\chapter{The Pauli basis test}
\label{chap:qpbt-test}
% Paper label sec:pauli-verifier kept for cross-reference fidelity; it resolves to this chapter.
\label{sec:pauli-verifier}

This chapter defines the two testing primitives underlying the question-reduction step: the \emph{classical low individual degree game}, in its simultaneous form, and the \emph{Pauli basis test}, a nonlocal game built from the classical game together with the Magic Square game. We state the quantum soundness of the classical game (Theorem~\ref{lem:ld-soundness}, imported, with the facts required by the identification, which remain to be proved, recorded in \cite{gap:qpbt_ld-dimension-divisibility}; see Remark~\ref{rem:ld-soundness-provider}), the rigidity of the Magic Square game (Theorem~\ref{thm:ms-rigidity}, imported), completeness of the Pauli basis test (Lemma~\ref{lem:pauli-completeness}), and the main soundness theorem for the Pauli basis test (Theorem~\ref{thm:pauli}), whose proof occupies Chapters~\ref{chap:qpbt-observables}--\ref{chap:qpbt-extraction} and is completed in Chapter~\ref{chap:qpbt-extraction} only modulo the proof obligation documented in \cite{gap:qpbt_ld-dimension-divisibility}. The chapter closes with the qubit form of the soundness statement (Corollary~\ref{cor:pauli-binary}) and the canonical parameter choice with its soundness bound (Definition~\ref{def:introparams}, Lemma~\ref{lem:delta-bound}).

Throughout, the two players are denoted $\mathrm A$ and $\mathrm B$; for $w \in \{\mathrm A,\mathrm B\}$ we write $\overline w$ for the other player. Games, strategies, values, conditionally linear (CL) functions and distributions, typed distributions sampled from the edges of a graph, and the state-dependent approximation $\approx_\delta$ between families of measurement operators are as in Chapter~\ref{chap:qpbt-games}; the algebraic material on lines, canonical linear maps, the low-degree encoding, and generalized Pauli observables is from Chapter~\ref{chap:qpbt-algebra}.

\section{The classical low individual degree game}
% Paper labels sec:ld-verifier and sec:ld-game kept for cross-reference fidelity.
\label{sec:ld-verifier}
\label{sec:ld-game}

The classical low individual degree game is a two-player game based on the low individual degree variant of the multilinearity test of Babai, Fortnow, and Lund. In the basic version, one player (the ``points player'') receives a uniformly random point $u \in \F_q^m$ and answers with a value $a \in \F_q$, while the other player (the ``lines player'') receives a random line through $u$ --- either axis-parallel or diagonal --- and answers with a univariate polynomial describing the restriction of a global function to that line. If the players share a polynomial $g\colon \F_q^m \to \F_q$ of individual degree at most $d$ and answer with $g(u)$ and with the restriction $g|_\ell$, they win with probability $1$; the soundness theorem below states that the converse approximately holds. In the simultaneous version, the players' answers are $k$-tuples, checked coordinatewise against a $k$-tuple of low individual degree polynomials.

\begin{definition}[The classical low individual degree game: parameters and registers]\label{def:ld-game}
	\lean{MIPStarRE.QPBT.ldGame}
	\leanok
	\uses{def:game, def:admissible-size}
	% The source fixes \N to be the set of positive integers
	% (references/qpbt-paper/04_preliminaries.tex:6) and takes m, d, ldc in \N
	% (references/qpbt-paper/08_classical_and_quantum_low_degree_tests.tex:33-35),
	% so the positivity below is the source's implicit domain, stated explicitly.
	% It is load-bearing: at d = 0 or k = 0 the prefactor a(dmk)^a of delta_ld in
	% lem:ld-soundness vanishes, and at d = 0 the resulting exact conclusions fail
	% (e.g. at (q,m,d,k) = (2,1,0,1) with eps = 1). The case k = 0, though it would
	% hold trivially, is likewise outside the source's domain and is omitted.
	The game $\mathfrak{G}^{\mathrm{ld}}$ is parametrized by a tuple $\mathsf{ldparams} = (q,m,d,k)$, where $m, d, k \ge 1$ are integers, $q \in \N$ is an admissible field size (Definition~\ref{def:admissible-size}), and $m$ divides $q$. We write $\mathfrak{G}^{\mathrm{ld}}_{\mathsf{ldparams}}$ to emphasize the dependence on the parameter tuple. The type set of $\mathfrak{G}^{\mathrm{ld}}$ is
	\[
		\mathcal T^{\mathrm{ld}} = \{\mathsf{Point}, \mathsf{ALine}, \mathsf{DLine}\}\;,
	\]
	and the distribution over type pairs $(t_{\mathrm A}, t_{\mathrm B})$ is uniform over $\mathcal T^{\mathrm{ld}} \times \mathcal T^{\mathrm{ld}}$. The ambient space is the direct sum $V = V_X \oplus V_I \oplus V_V$ of three register subspaces: the \emph{point register} $V_X \cong \F_q^m$, the \emph{coordinate register} $V_I \cong \F_q$, and the \emph{direction register} $V_V \cong \F_q^m$. Elements of $V$ are identified with triples $(u,s,v) \in \F_q^m \times \F_q \times \F_q^m$. The question distribution and the win predicate of $\mathfrak{G}^{\mathrm{ld}}$ are specified in Definitions~\ref{def:ld-question-distribution} and~\ref{def:ld-win-predicate} below.
\end{definition}

\begin{definition}[Question distribution of the classical low individual degree game]\label{def:ld-question-distribution}
	\uses{def:ld-game, def:cl-func, def:cl-dist, lem:cl-concat, def:cl-canonical, def:line-representative}
	Let $\mathsf{ldparams} = (q,m,d,k)$ be as in Definition~\ref{def:ld-game}. Define $L_{\mathsf{Point}}$ to be the $1$-level CL function on $V$ that projects onto $V_X$: for every $(u,s,v) \in V$,
	\begin{equation}
		\label{eq:cl-ptf}
		L_{\mathsf{Point}}(u,s,v) = (u,0,0)\;.
	\end{equation}
	For $s \in \F_q$, let $\chi(s)$ be the unique integer such that
	\begin{equation}
		\label{eq:chi-func}
		s = (\chi(s)-1)\,\frac{q}{m} + r
	\end{equation}
	for some $0 \le r < q/m$, where $s$ is interpreted as an integer in $\{0,1,\ldots,q-1\}$ under the identification of $\F_q$ with binary strings fixed in Chapter~\ref{chap:qpbt-algebra}. Since $m$ divides $q$, when $s$ is uniformly distributed on $\F_q$ the index $\chi(s)$ is uniformly distributed on $\{1,2,\ldots,m\}$. Define $L_{\mathsf{ALine}}$ by
	\begin{equation}
		\label{eq:cl-alnf}
		L_{\mathsf{ALine}}(u,s,v) = \bigl(L^{\mathrm{Ln}}_{e_i}(u),\, s,\, 0\bigr)\quad\text{for all } (u,s,v) \in V\;,
	\end{equation}
	where $i = \chi(s)$ and $L^{\mathrm{Ln}}_v$ denotes the canonical linear map with one-dimensional kernel spanned by $v$, used to compute canonical representatives of lines (Definition~\ref{def:line-representative}). The function $L_{\mathsf{ALine}}$ is the concatenation (in the sense of Lemma~\ref{lem:cl-concat}) of the $1$-level CL function projecting $V_I \oplus V_V$ onto $V_I$ with the family of $1$-level CL functions $\{L^{\mathrm{Ln}}_{e_i}\}$ on $V_X$ indexed by $i \in \{1,\ldots,m\}$; thus $L_{\mathsf{ALine}}$ is a $2$-level CL function. Define $L_{\mathsf{DLine}}$ by
	\begin{equation}
		\label{eq:cl-dlnf}
		L_{\mathsf{DLine}}(u,s,v) = \bigl(L^{\mathrm{Ln}}_{\pi_{i-1}(v)}(u),\, s,\, \pi_{i-1}(v)\bigr)\;,
	\end{equation}
	where $i = \chi(s)$ and $\pi_{i-1}$ zeroes out the first $i-1$ coordinates of its input. The function $L_{\mathsf{DLine}}$ is the concatenation of (i) the identity on $V_I$, (ii) for each $i$ the $1$-level CL function $v \mapsto \pi_{i-1}(v)$ on $V_V$, and (iii) for each pair $(s,v') \in V_I \oplus V_V$ the $1$-level CL function $L^{\mathrm{Ln}}_{v'}$ on $V_X$; thus $L_{\mathsf{DLine}}$ is a $3$-level CL function.

	The question distribution of $\mathfrak{G}^{\mathrm{ld}}$ is the typed CL distribution obtained as follows: sample a type pair $(t_{\mathrm A}, t_{\mathrm B})$ uniformly from $\mathcal T^{\mathrm{ld}} \times \mathcal T^{\mathrm{ld}}$, sample $z = (u,s,v)$ uniformly from $V$, and send each player $w \in \{\mathrm A,\mathrm B\}$ the question $(t_w, L_{t_w}(z))$.
\end{definition}

\begin{lemma}[Level of the point map]\label{lem:ld-point-level}
	\lean{MIPStarRE.QPBT.isCondLinear_ldPointCL}
	\leanok
	\uses{def:ld-question-distribution, def:cl-func}
	Fix $\mathsf{ldparams} = (q,m,d,k)$ as in Definition~\ref{def:ld-game}. The map $L_{\mathsf{Point}}$ of Equation~\eqref{eq:cl-ptf} is a $1$-level conditionally linear function on $V$.
\end{lemma}
\begin{proof}
	\leanok
	\uses{def:ld-question-distribution, def:cl-func}
	The map $L_{\mathsf{Point}}$ is the linear projection onto the point register $V_X$, hence is a $1$-level CL function.
\end{proof}

\begin{lemma}[Level of the axis-parallel line map]\label{lem:ld-aline-level}
	\lean{MIPStarRE.QPBT.isCondLinear_ldALineCL}
	\leanok
	\uses{def:ld-question-distribution, def:cl-func}
	Fix $\mathsf{ldparams} = (q,m,d,k)$ as in Definition~\ref{def:ld-game}. The map $L_{\mathsf{ALine}}$ of Equation~\eqref{eq:cl-alnf} is a $2$-level conditionally linear function on $V$.
\end{lemma}
\begin{proof}
	\leanok
	\uses{def:cl-func, def:register-subspace, def:ld-question-distribution, def:line-representative}
	The following data exhibit $L_{\mathsf{ALine}}$ as a $2$-level conditionally linear function. Its first level is the projection of $V$ onto the coordinate register $V_I$, which is linear. For a first-level value with coordinate component $s$, its second level is the linear map $L^{\mathrm{Ln}}_{e_i}$ with $i = \chi(s)$, acting on the point register $V_X$. The registers $V_I$ and $V_X$ are disjoint, so the second level is supported on a register subspace disjoint from the first, and summing the two contributions gives $\bigl(L^{\mathrm{Ln}}_{e_i}(u),\, s,\, 0\bigr)$, which is $L_{\mathsf{ALine}}(u,s,v)$ by Equation~\eqref{eq:cl-alnf}.
\end{proof}

\begin{lemma}[Level of the diagonal line map]\label{lem:ld-dline-level}
	\lean{MIPStarRE.QPBT.isCondLinear_ldDLineCL}
	\leanok
	\uses{def:ld-question-distribution, def:cl-func}
	Fix $\mathsf{ldparams} = (q,m,d,k)$ as in Definition~\ref{def:ld-game}. The map $L_{\mathsf{DLine}}$ of Equation~\eqref{eq:cl-dlnf} is a $3$-level conditionally linear function on $V$.
\end{lemma}
\begin{proof}
	\leanok
	\uses{def:cl-func, def:register-subspace, def:ld-question-distribution, def:line-representative}
	The following data exhibit $L_{\mathsf{DLine}}$ as a $3$-level conditionally linear function. Its first level is the projection of $V$ onto the coordinate register $V_I$. For a first-level value with coordinate component $s$, its second level is the linear map $v \mapsto \pi_{i-1}(v)$ with $i = \chi(s)$ on the direction register $V_V$. For a second-level value with direction component $v'$, its third level is the linear map $L^{\mathrm{Ln}}_{v'}$ on the point register $V_X$. The registers $V_I$, $V_V$ and $V_X$ are pairwise disjoint, so each level is supported on a register subspace disjoint from the preceding ones, and summing the three contributions gives $\bigl(L^{\mathrm{Ln}}_{\pi_{i-1}(v)}(u),\, s,\, \pi_{i-1}(v)\bigr)$, which is $L_{\mathsf{DLine}}(u,s,v)$ by Equation~\eqref{eq:cl-dlnf}.
\end{proof}

\begin{lemma}[Formalization support for the low-degree typed CL family]\label{lem:ld-question-cl-family}
	\lean{MIPStarRE.QPBT.isTypedCondLinearFamily_ldCL}
	\leanok
	\uses{def:ld-question-distribution, def:typed-cl-functions}
	This is not a separately named statement of the source article; it isolates the family condition in Lemma~\ref{lem:ld-question-typed-cl}. For every low-degree type $t \in \mathcal T^{\mathrm{ld}}$, the map $L_t$ of Definition~\ref{def:ld-question-distribution} is a $3$-level CL function. Thus $L = \{L_t\}_{t \in \mathcal T^{\mathrm{ld}}}$ is a finite typed family of $3$-level CL functions.
\end{lemma}
\begin{proof}
	\leanok
	\uses{lem:cl-level-mono, lem:ld-point-level, lem:ld-aline-level, lem:ld-dline-level}
	The point, axis-line, and diagonal-line maps have levels $1$, $2$, and $3$, respectively. Raise the first two representations to level $3$ by Lemma~\ref{lem:cl-level-mono}.
\end{proof}

\begin{lemma}[Formalization support for the low-degree typed distribution equality]\label{lem:ld-question-typed-equality}
	\lean{MIPStarRE.QPBT.ldQuestionDistribution_eq_typedCL}
	\leanok
	\uses{def:ld-question-distribution, def:typed-cl-distributions, def:graph-distribution}
	This is not a separately named statement of the source article; it isolates the distribution equality in Lemma~\ref{lem:ld-question-typed-cl}. Fix $\mathsf{ldparams} = (q,m,d,k)$ as in Definition~\ref{def:ld-game}, let $L = \{L_t\}_{t \in \mathcal T^{\mathrm{ld}}}$ be the family of Definition~\ref{def:ld-question-distribution}, and let $G^{\mathrm{ld}}$ be the complete type graph. The question distribution of $\mathfrak G^{\mathrm{ld}}$ equals the distribution $\mu^{G^{\mathrm{ld}}}_{L,L}$ that the construction of Definition~\ref{def:typed-cl-distributions} produces from $(L,L)$ and $G^{\mathrm{ld}}$. This statement records only the distribution identity; Lemma~\ref{lem:ld-question-cl-family} supplies the common-level typed-family assertion.
\end{lemma}
\begin{proof}
	\leanok
	\uses{def:ld-question-distribution, def:typed-cl-distributions, def:graph-distribution, def:cl-dist, lem:distribution-map-comp, lem:uniform-bind-map}
	Every unordered pair of types is an edge of $G^{\mathrm{ld}}$, so the graph distribution $\mu_{G^{\mathrm{ld}}}$ of Definition~\ref{def:graph-distribution} is uniform on $\mathcal T^{\mathrm{ld}} \times \mathcal T^{\mathrm{ld}}$. For each ordered type pair $(t_{\mathrm A}, t_{\mathrm B})$, the CL distribution $\mu_{L_{t_{\mathrm A}},\, L_{t_{\mathrm B}}}$ is the push-forward of the uniform distribution on $V$ along $z \mapsto \bigl(L_{t_{\mathrm A}}(z), L_{t_{\mathrm B}}(z)\bigr)$. Binding the uniform type pair to these push-forwards and tagging each content with its type is exactly the two-stage sampling procedure of Definition~\ref{def:ld-question-distribution}.
\end{proof}

\begin{lemma}[The low-degree question distribution is a typed CL distribution]\label{lem:ld-question-typed-cl}
	\lean{MIPStarRE.QPBT.ldQuestionDistribution_isTypedCL}
	\leanok
	\uses{def:ld-question-distribution, def:typed-cl-distributions}
	Fix $\mathsf{ldparams} = (q,m,d,k)$ as in Definition~\ref{def:ld-game}, let $L = \{L_t\}_{t \in \mathcal T^{\mathrm{ld}}}$ be the family of Definition~\ref{def:ld-question-distribution}, and let $G^{\mathrm{ld}}$ be the graph on $\mathcal T^{\mathrm{ld}}$ all of whose unordered pairs, self-loops included, are edges. The family $L$ is a finite $\mathcal T^{\mathrm{ld}}$-typed family of $3$-level conditionally linear functions, and the question distribution of $\mathfrak{G}^{\mathrm{ld}}$ equals the $(\mathcal T^{\mathrm{ld}}, G^{\mathrm{ld}})$-typed conditionally linear distribution $\mu^{G^{\mathrm{ld}}}_{L,\,L}$ of Definition~\ref{def:typed-cl-distributions}.
\end{lemma}
\begin{proof}
	\leanok
	\uses{lem:ld-question-cl-family, lem:ld-question-typed-equality}
	Combine Lemma~\ref{lem:ld-question-cl-family}, which supplies the typed-family hypothesis, with Lemma~\ref{lem:ld-question-typed-equality}, which identifies the two distributions.
\end{proof}

The next two lemmas show that the CL functions $L_{\mathsf{Point}}$, $L_{\mathsf{ALine}}$, $L_{\mathsf{DLine}}$ give rise to the axis-parallel and diagonal line-versus-point distributions.
We label the coordinates by $r\in\{0,\ldots,m-1\}$, with label $r$
representing the paper's coordinate $r+1$.

\begin{lemma}\label{lem:alnf}
	\lean{MIPStarRE.QPBT.aLinePointDist_point_marginal_uniform, MIPStarRE.QPBT.aLinePointDist_mem_line}
	\leanok
	\uses{def:line-point-dist, def:line}
	Consider the pair $(\ell, u)$ sampled in the following way: first, sample the tuple $\bigl((u_0,s,0),(u,0,0)\bigr)$ from the CL distribution $\mu_{L_{\mathsf{ALine}}, L_{\mathsf{Point}}}$, and let $\ell = \ell(u_0, e_i)$ where $i = \chi(s)$. Then $u$ is a uniformly random point in $\F_q^m$, and $\ell$ is a line that goes through $u$ and is parallel to a uniformly random axis $i \in \{1,2,\ldots,m\}$.
\end{lemma}
\begin{proof}
	\leanok
	\uses{def:cl-dist, def:ld-question-distribution, def:ld-line-description, def:line-representative, lem:distribution-map-comp, lem:uniform-map-equivalence, lem:uniform-product-first-marginal, lem:uniform-product-second-marginal, lem:chi-index-uniform, lem:line-representative-idempotent, lem:line-rep-incidence}
	By Definition~\ref{def:cl-dist} the distribution $\mu_{L_{\mathsf{ALine}}, L_{\mathsf{Point}}}$ is the push-forward of the uniform distribution on $V$ along $z \mapsto \bigl(L_{\mathsf{ALine}}(z), L_{\mathsf{Point}}(z)\bigr)$, so each marginal below is the push-forward of the uniform distribution on $V$ along the composite of that map with the corresponding coordinate. The point marginal is the image of the uniform distribution on $V$ under $(u,s,v) \mapsto u$, hence uniform on $\F_q^m$. The axis marginal is its image under $(u,s,v) \mapsto \chi(s)$; since $m$ divides $q$, every fibre of $\chi$ has exactly $q/m$ elements, so $i = \chi(s)$ is uniform on $\{1,2,\ldots,m\}$ and $(u_0,s)$ specifies a uniformly random axis-parallel line $\ell(u_0,e_i)$. Finally the base point recorded in the sample is $u_0 = L^{\mathrm{Ln}}_{e_i}(u)$, which is fixed by $L^{\mathrm{Ln}}_{e_i}$ (Definition~\ref{def:line-representative}) and satisfies $u \in \ell(u_0,e_i)$ by Lemma~\ref{lem:line-rep-incidence}.
\end{proof}

\begin{lemma}\label{lem:dlnf}
	\lean{MIPStarRE.QPBT.dLinePointDist_point_marginal_uniform, MIPStarRE.QPBT.dLinePointDist_mem_line, MIPStarRE.QPBT.dLinePointDist_prefix_zero}
	\leanok
	\uses{def:line-point-dist, def:line}
	Consider the pair $(\ell, u)$ sampled in the following way: first, sample the tuple $\bigl((u_0,s,v),(u,0,0)\bigr)$ from the CL distribution $\mu_{L_{\mathsf{DLine}}, L_{\mathsf{Point}}}$, and let $\ell = \ell(u_0, v)$ where $i = \chi(s)$. Then $u$ is a uniformly random point in $\F_q^m$, and $\ell$ is a random diagonal line that goes through $u$ and shares its first $i-1$ coordinates with $u$, for uniformly random $i \in \{1,2,\ldots,m\}$.
\end{lemma}
\begin{proof}
	\leanok
	\uses{def:cl-dist, def:ld-question-distribution, def:ld-line-description, def:line-representative, lem:distribution-map-comp, lem:uniform-map-equivalence, lem:uniform-product-first-marginal, lem:uniform-product-second-marginal, lem:chi-index-uniform, lem:prefix-projection-idempotent, lem:line-representative-idempotent, lem:line-description-prefix, lem:line-rep-incidence}
	The three steps are those of the proof of Lemma~\ref{lem:alnf}, and none of them depends on that lemma. The two marginals are again push-forwards of the uniform distribution on $V$ along coordinates of $z \mapsto \bigl(L_{\mathsf{DLine}}(z), L_{\mathsf{Point}}(z)\bigr)$, by Definition~\ref{def:cl-dist}: the point marginal is uniform on $\F_q^m$, and $i = \chi(s)$ is uniform on $\{1,2,\ldots,m\}$ because every fibre of $\chi$ has $q/m$ elements. The direction component of the sampled tuple is $v' = \pi_{i-1}(v)$, whose first $i-1$ coordinates vanish and which is unchanged by a further application of $\pi_{i-1}$. Finally $u_0 = L^{\mathrm{Ln}}_{v'}(u)$ is fixed by $L^{\mathrm{Ln}}_{v'}$ (Definition~\ref{def:line-representative}) and satisfies $u \in \ell(u_0,v')$ by Lemma~\ref{lem:line-rep-incidence}.
\end{proof}

\begin{definition}[Line-point distributions]\label{def:line-point-dist}
	% A line description retains the seed used to select its coordinate.
	\lean{MIPStarRE.QPBT.aLinePointDist, MIPStarRE.QPBT.dLinePointDist}
	\lean{MIPStarRE.QPBT.linePointDist}
	\lean{MIPStarRE.QPBT.aLinePointDist_isProbability}
	\lean{MIPStarRE.QPBT.dLinePointDist_isProbability}
	\lean{MIPStarRE.QPBT.linePointDist_isProbability}
	% The refined laws retain data not present in the axis source law.
	\uses{def:ld-question-distribution, def:ld-line-description}
	The source \emph{axis line-point distribution} $D_{\mathsf{ALine}}$ is the law $\mu_{L_{\mathsf{ALine}}, L_{\mathsf{Point}}}$ on pairs $(\ell=(u_0,e_i),u)$, and the source \emph{diagonal line-point distribution} $D_{\mathsf{DLine}}$ is the corresponding law on pairs $(\ell=(u_0,s,v),u)$ from $\mu_{L_{\mathsf{DLine}}, L_{\mathsf{Point}}}$. Their equal mixture is $D_{\mathsf{Line}}$. We also use seed-bearing refinements of these laws. An axis description records $u_0$ and a seed $s$ selecting $i=\chi(s)$, subject to $L^{\mathrm{Ln}}_{e_i}(u_0)=u_0$. A diagonal description records $(u_0,s,v)$, subject to $L^{\mathrm{Ln}}_v(u_0)=u_0$ and to the prefix condition $v_j=0$ for $j<i$. Mapping the two kinds of description to $(u_0,e_{\chi(s)})$ and $(u_0,s,v)$, respectively, sends the refined laws to the source laws.
\end{definition}

\begin{definition}[Win predicate of the classical low individual degree game]\label{def:ld-win-predicate}
	\uses{def:ld-game, def:ld-question-distribution, def:line, def:polyfunc}
	Fix $\mathsf{ldparams} = (q,m,d,k)$. Answers in the game $\mathfrak{G}^{\mathrm{ld}}_{\mathsf{ldparams}}$ take the following forms, depending on the question type; answers not of the prescribed form are rejected.
	\begin{itemize}
		\item Type $\mathsf{Point}$, question content $u \in \F_q^m$: the answer is an element of $\F_q^k$.
		\item Type $\mathsf{ALine}$, question content $(u_0,s) \in \F_q^m \times \F_q$, specifying the axis-parallel line $\ell(u_0,e_i)$ with $i = \chi(s)$: the answer is a $k$-tuple $(f_1,\ldots,f_k)$ of univariate polynomials $f_j\colon \F_q \to \F_q$ of degree at most $d$, each given by its $d+1$ coefficients.
		\item Type $\mathsf{DLine}$, question content $(u_0,s,v) \in \F_q^m \times \F_q \times \F_q^m$, specifying the diagonal line $\ell(u_0,v')$ with $v' = \pi_{i-1}(v)$ and $i = \chi(s)$: the answer is a $k$-tuple of univariate polynomials of degree at most $md$.
	\end{itemize}
	The win predicate $D^{\mathrm{ld}} = D^{\mathrm{ld}}_{\mathsf{ldparams}}$ is the function of the tuple $(t_{\mathrm A}, x_{\mathrm A}, t_{\mathrm B}, x_{\mathrm B}, a_{\mathrm A}, a_{\mathrm B})$ determined by the following clauses, evaluated for $w \in \{\mathrm A,\mathrm B\}$; in all cases where no action is indicated, the predicate accepts.
	\begin{enumerate}
		\item (\textbf{Consistency test}) If $t_{\mathrm A} = t_{\mathrm B}$, accept if and only if $a_{\mathrm A} = a_{\mathrm B}$.
		% The source phrases both clauses with "the" t such that x = u_0 + tv; that phrase
		% does not determine t on diagonal questions with direction v' = 0, which the
		% question distribution samples with positive probability. The universal
		% quantification over t used here is the local correction; see the remark below.
		\item (\textbf{Axis-parallel line-versus-point test}) If $t_w = \mathsf{ALine}$ and $t_{\overline w} = \mathsf{Point}$, accept if and only if $f_j(t) = (a_{\overline w})_j$ for all $j \in \{1,2,\ldots,k\}$ and all $t \in \F_q$ such that $x_{\overline w} = u_0 + t e_i$, where $i = \chi(s)$.
		\item (\textbf{Diagonal line-versus-point test}) If $t_w = \mathsf{DLine}$ and $t_{\overline w} = \mathsf{Point}$, accept if and only if $f_j(t) = (a_{\overline w})_j$ for all $j \in \{1,2,\ldots,k\}$ and all $t \in \F_q$ such that $x_{\overline w} = u_0 + t v'$, where $v' = \pi_{i-1}(v)$ and $i = \chi(s)$.
	\end{enumerate}
	The game $\mathfrak{G}^{\mathrm{ld}}_{\mathsf{ldparams}}$ is the typed game with type set $\mathcal T^{\mathrm{ld}}$, the question distribution of Definition~\ref{def:ld-question-distribution}, and win predicate $D^{\mathrm{ld}}_{\mathsf{ldparams}}$.
\end{definition}

\begin{remark}[Diagonal questions with zero direction]\label{rem:ld-win-zero-direction}
	The source states the two line-versus-point clauses with ``the'' $t \in \F_q$ such that $x_{\overline w} = u_0 + t v$. When the direction is nonzero and $x_{\overline w}$ lies on the line, exactly one such $t$ exists and the formulation above agrees with the source. But zero directions occur: Definition~\ref{def:line} explicitly permits them, and the question distribution of Definition~\ref{def:ld-question-distribution} produces $\mathsf{DLine}$ questions with $v' = \pi_{i-1}(v) = 0$ with positive probability, namely whenever the last $m - i + 1$ coordinates of $v$ vanish. In that case $\ell(u_0, 0) = \{u_0\}$, every $t \in \F_q$ satisfies $x_{\overline w} = u_0 + t v'$ when $x_{\overline w} = u_0$, and the source clause does not determine acceptance for a nonconstant answer polynomial. The universal quantification over $t$ adopted in Definition~\ref{def:ld-win-predicate} makes the predicate a well-defined function of its arguments; it accepts vacuously when no such $t$ exists, a case the question distribution never produces. Completeness is unaffected: the restriction of a polynomial $g$ to the singleton line $\ell(u_0,0)$, expressed through the parameterization $t \mapsto u_0 + t \cdot 0$, is the constant polynomial with value $g(u_0)$, which passes the check for every $t$. Soundness is likewise unaffected: zero directions occur with probability at most $2/(mq)$, so any acceptance convention on this event changes success probabilities by at most that amount; the game-correspondence obligation enumerated in the proof of Theorem~\ref{lem:ld-soundness} accounts for this event explicitly, and an error of this shape is absorbed into the universal constants of $\delta_{\mathrm{ld}}$.
\end{remark}

The soundness properties of the low individual degree game are expressed in terms of the following class of measurements.

\begin{definition}[Low-degree polynomial measurements]\label{def:ld-meas}
	\lean{MIPStarRE.QPBT.PolyIndex, MIPStarRE.QPBT.PolyMeas, MIPStarRE.QPBT.PolyMeasFamily, MIPStarRE.QPBT.PolyMeasTuple}
	\leanok
	\uses{def:polyfunc, def:polymeasurements, def:polynomials-degree}
	Let $\polymeas{m}{q}{d}$ denote the set of POVM measurements indexed by $\polyfunc{m}{q}{d}$, the polynomial representatives of individual degree at most $d$ in each variable; this is the notion of Definition~\ref{def:polymeasurements}, and the argument order $(m,q,d)$ follows that definition (the paper orders the arguments as $(m,d,q)$). By Definition~\ref{def:polynomials-degree}, evaluation identifies these representatives with the corresponding polynomial functions $g\colon \F_q^m \to \F_q$, bijectively whenever $d \le q-1$; at the canonical parameters (Definition~\ref{def:introparams}, where $d = 1 \le q-1$) the identification is bijective; the general statements of this chapter quantify over all admissible tuples and are read through the representative convention. More generally, for an integer $k$ and tuples $m = (m_1,\ldots,m_k)$, $q = (q_1,\ldots,q_k)$, $d = (d_1,\ldots,d_k)$, let $\mathrm{PolyMeas}(m,q,d,k)$ denote the set of measurements $G = \{G_{g_1,\ldots,g_k}\}$ such that for each $i \in \{1,\ldots,k\}$, the outcome index $g_i$ is a polynomial $g_i\colon \F_{q_i}^{m_i} \to \F_{q_i}$ of individual degree at most $d_i$.
	The tuple components are labelled by $r\in\{0,\ldots,k-1\}$, with
	label $r$ representing the paper's component $r+1$.
\end{definition}

Quantum soundness of a version of the classical low-degree test for polynomials of low \emph{total} degree was claimed by Natarajan and Vidick, building on an analysis of a three-prover version of the test by Vidick; there is a known gap in the soundness analysis of that test, and it is not used here. Instead, the construction relies on the low \emph{individual} degree test, whose quantum soundness is proved in~\cite{Ji2020LowIndividualDegree} and generalized in~\cite{Ji2021Quantum}. The statement below is the form most directly useful in what follows.

% The paper states this result in a theorem environment with the label lem:ld-soundness.
% We keep both the environment and the label verbatim for cross-reference fidelity.
\begin{theorem}[Quantum soundness of the simultaneous classical low individual degree test]\label{lem:ld-soundness}
	\lean{MIPStarRE.QPBT.deltaLd, MIPStarRE.QPBT.exists_ld_soundness}
	\leanok
	\uses{def:ld-game, def:ld-question-distribution, def:ld-win-predicate, def:ld-meas, def:projective-strategy-general, def:tensor-product-value, def:consistency, def:bracket}
	There exists a function
	\[
		\delta_{\mathrm{ld}}(\eps, q, m, d, k) = a\,(dmk)^a \bigl(\eps^b + q^{-b} + 2^{-bmd}\bigr)
	\]
	for universal constants $a \ge 1$ and $0 < b \le 1$ such that the following holds. For all $\eps > 0$ and every parameter tuple $\mathsf{ldparams} = (q,m,d,k)$ as in Definition~\ref{def:ld-game} (so $m, d, k \ge 1$, $q$ is an admissible field size, and $m$ divides $q$), and for all projective strategies $(\psi, A, B)$ that succeed with probability at least $1-\eps$ in the game $\mathfrak{G}^{\mathrm{ld}}_{\mathsf{ldparams}}$, there exist measurements
	\[
		G^w \in \mathrm{PolyMeas}(m,q,d,k)
	\]
	on $\hilb_w$ for $w \in \{\mathrm A,\mathrm B\}$, where by slight abuse of notation $m$, $q$, $d$ denote the constant tuples $(m,\ldots,m)$, $(q,\ldots,q)$, $(d,\ldots,d)$ of length $k$, such that
	\begin{align*}
		A^{\mathsf{Point},\, u}_{a_1,\ldots,a_k} \ot I_{\mathrm B}
		&\simeq_{\delta_{\mathrm{ld}}}
		I_{\mathrm A} \ot G^{\mathrm B}_{[(g_1(u),\ldots,g_k(u)) = (a_1,\ldots,a_k)]}\;,\\
		G^{\mathrm A}_{[(g_1(u),\ldots,g_k(u)) = (a_1,\ldots,a_k)]} \ot I_{\mathrm B}
		&\simeq_{\delta_{\mathrm{ld}}}
		I_{\mathrm A} \ot B^{\mathsf{Point},\, u}_{a_1,\ldots,a_k}\;,\\
		G^{\mathrm A}_{g_1,\ldots,g_k} \ot I_{\mathrm B}
		&\simeq_{\delta_{\mathrm{ld}}}
		I_{\mathrm A} \ot G^{\mathrm B}_{g_1,\ldots,g_k}\;,
	\end{align*}
	where $\delta_{\mathrm{ld}} = \delta_{\mathrm{ld}}(\eps,q,m,d,k)$, all three relations are the state-dependent consistency $\simeq$ of Definition~\ref{def:consistency} with respect to the state $\ket{\psi}$, and the bracket notation for post-processed outcomes is that of Definition~\ref{def:bracket}. The first two relations hold on average over the uniform distribution on $u \in \F_q^m$; no question distribution is associated with the third.
\end{theorem}
\begin{proof}
	\uses{def:ld-question-distribution, def:ld-win-predicate, def:ld-meas, def:povm-distance}
	Let $\Sigma = \F_q$ and let $\code \subseteq \Sigma^q$ be the Reed--Solomon code of degree $d$: the set of all $(p(x_1),\ldots,p(x_q))$ where $p\colon \F_q \to \F_q$ is a univariate polynomial of degree at most $d$ and $x_1,\ldots,x_q$ enumerate $\F_q$. We may assume $d < q$. In the opposite regime the statement is immediate: for $d \ge q$ the prefactor gives $\delta_{\mathrm{ld}} \ge a\,(dmk)^a\, q^{-b} \ge q^{a-b} \ge 1$, using $m, k \ge 1$, while the quantity bounded in Definition~\ref{def:consistency} never exceeds $4$ for families of POVM elements: $\sum_a \lVert (A_a \ot I - I \ot B_a)\ket{\psi}\rVert^2 \le 2\sum_a \bigl(\bra{\psi} A_a^2 \ot I \ket{\psi} + \bra{\psi} I \ot B_a^2 \ket{\psi}\bigr) \le 4$, using $0 \le M_a^2 \le M_a$ and $\sum_a M_a = \Id$ for POVM elements, so the three relations hold for an arbitrary choice of $G^w \in \mathrm{PolyMeas}(m,q,d,k)$. For $d < q$, evaluation is injective on polynomials of degree at most $d$ and a nonzero polynomial of degree at most $d$ vanishes at no more than $d$ of the $q$ points, so $\code$ is a linear $[q,\, d+1,\, q-d]_\Sigma$ code.
	% The paper prints the parameters as [q, d+1, n-d+1] with an undefined n. With the
	% intended value n = q the printed distance q-d+1 exceeds the true minimum distance
	% q-d of the degree-d evaluation code, and the dimension claim d+1 requires d < q.
	% Both points are corrected above; the regime d >= q, which the parameter tuple of
	% Definition def:ld-game permits, is handled by the triviality argument. The
	% boundary values d = 0 and k = 0 are excluded by the parameter domain of
	% Definition def:ld-game, following the source's convention that \N denotes
	% the positive integers.
	% The source asserts a flat identity between its game and the two-prover
	% tensor-code test and applies Theorem 4.7 of ji2021quantum with the
	% parameter choice m^3 d
	% (references/qpbt-paper/08_classical_and_quantum_low_degree_tests.tex:446-454),
	% without checking that theorem's hypothesis k >= 12mt on the parameter.
	% Inspection of the cited test (arXiv:2111.08131, Figures 1 and 2,
	% Definitions 4.5-4.6, Theorem 4.7) shows the asserted identity is not one:
	% the question formats, the branch weights, and the zero-direction
	% convention all differ, and the parameter bound fails at small m. Per the
	% import policy, this proof states the source's route and enumerates
	% the facts required by the identification, which remain to be proved, as open obligations; the
	% detailed comparison and the case arithmetic are in the import-obligations
	% section of docs/paper-gaps/qpbt_ld-dimension-divisibility.tex.
	In this regime the statement is imported, and the source's route is as follows: identify the game $\mathfrak{G}^{\mathrm{ld}}_{\mathsf{ldparams}}$ with $k = 1$ with the two-prover tensor-code test for the code $\code^{\ot m}$ of~\cite{Ji2021Quantum} (Figures~1 and~2 there, the roles of the two players being assigned uniformly at random); apply Theorem~4.7 of~\cite{Ji2021Quantum}, instantiated with its auxiliary integer parameter --- denoted $k$ there and unrelated to the simultaneity parameter $k$ of $\mathsf{ldparams}$ --- chosen as $m^3 d$, so that its error function $\mathrm{poly}(m, d+1, m^3d) \cdot \mathrm{poly}(\eps, q^{-1}, e^{-md})$ takes the displayed form of $\delta_{\mathrm{ld}}$; and extend from $k = 1$ to general $k$ by a standard reduction, following exactly the derivation of Theorem~4.43 from Theorem~4.40 in~\cite{Natarajan2018NEEXP}. The import carries two verification obligations, which remain open; they are enumerated here as obligations and detailed, together with candidate repairs, in the import-obligations section of~\cite{gap:qpbt_ld-dimension-divisibility}.
	\begin{enumerate}
		\item (\emph{Game correspondence.}) The source asserts that the two games are identical; they are not. In the cited test, the subcube commutation and pair-synchronicity branches send an ordered pair of points of a \emph{hidden} subcube and receive a pair of field values --- a question type that $\mathfrak{G}^{\mathrm{ld}}$ does not have --- its line questions carry no analogue of the seed $s$, its nontrivial cross-tests have mass $1/3$ where the present ones have $2/9$, and it fixes no acceptance convention on the zero-direction event of Remark~\ref{rem:ld-win-zero-direction}, whose probability is at most $\frac1m \sum_{j=1}^m q^{-j} \le \frac{1}{m(q-1)} \le \frac{2}{mq}$ over uniform $(s,v)$. What must be proved is a reduction converting any projective strategy for $\mathfrak{G}^{\mathrm{ld}}$ with success probability $1-\eps$ into a two-prover strategy for the tensor-code test with rejection probability at most $\tfrac32\,\eps + O(1/(mq))$ --- the factor $\tfrac32$ being the target bookkeeping constant from the branch weights, and the $O(1/(mq))$ the zero-direction term. A deficit of this shape is then absorbed into $\delta_{\mathrm{ld}}$ by enlarging the universal constant $a$, since $(x+y)^b \le x^b + y^b$ for $0 < b \le 1$ and $(2/(mq))^b \le 2\,q^{-b}$ for $m \ge 1$, using $dmk \ge 1$.
		\item (\emph{Parameter bound.}) Theorem~4.7 requires its auxiliary parameter to be at least $12mt$, where $t$ is the interpolation parameter of the base code, equal to $d+1$ for the degree-$d$ Reed--Solomon code $\code$ (blocklength $q$, distance $q-d$). The choice $m^3 d$ satisfies $m^3 d \ge 12m(d+1)$ if and only if $m^2 d \ge 12(d+1)$, which holds for all $d \ge 1$ when $m \ge 5$ and, at $m = 4$, exactly when $d \ge 3$, but fails for every $d$ when $m \le 3$. Over the parameter domain of Definition~\ref{def:ld-game}, where $m$ divides the admissible field size $q$ and is therefore a power of $2$, the bound thus fails exactly at $m \in \{1,2\}$ for every $d$ and at $m = 4$ for $d \in \{1,2\}$; the case check and the candidate repair --- instantiating the parameter as $\max(m^3 d,\, 12m(d+1))$ and re-absorbing the error terms --- are in the note.
	\end{enumerate}
	Quantum soundness of the low individual degree test itself is the result of~\cite{Ji2020LowIndividualDegree}; Remark~\ref{rem:ld-soundness-provider} relates the present statement to Theorem~\ref{thm:main-formal}, which proves that result.
\end{proof}

\begin{remark}[Relation to the low individual degree test of Chapter~\ref{chap:test}]\label{rem:ld-soundness-provider}
	Theorem~\ref{lem:ld-soundness} is imported: the proof above records the source's reduction to Theorem~4.7 of~\cite{Ji2021Quantum}, instantiated with that theorem's auxiliary parameter chosen as $m^3 d$ and followed by the extension to simultaneous $k$-tuples of~\cite{Natarajan2018NEEXP}, together with the two open facts required by the identification, which remain to be proved, --- the game correspondence and the parameter bound, detailed in~\cite{gap:qpbt_ld-dimension-divisibility}. The result of~\cite{Ji2020LowIndividualDegree} that underlies it is proved here as Theorem~\ref{thm:main-formal}: for every projective strategy passing the $(m,q,d)$ low individual degree test of Definition~\ref{def:low-individual-degree-test} with probability at least $1-\eps$, that theorem produces measurements in $\polymeas{m}{q}{d}$ satisfying consistency relations of the same shape as the three approximations of Theorem~\ref{lem:ld-soundness}. Theorem~\ref{lem:ld-soundness} is not a consequence of Theorem~\ref{thm:main-formal} alone: a derivation from it would require, in addition, an equivalence between the typed game $\mathfrak{G}^{\mathrm{ld}}_{(q,m,d,1)}$ and the $(m,q,d)$ test of Chapter~\ref{chap:test} --- matching the two question distributions, the two win conditions, and the two error functions --- together with the extension from $k = 1$ to general $k$ by the reduction of~\cite{Natarajan2018NEEXP} cited above. Neither implication is proved in this chapter, and the proof of Theorem~\ref{lem:ld-soundness} given above does not invoke Theorem~\ref{thm:main-formal}.
\end{remark}

% This remark records the formal status of the extension to general k;
% it is not a statement from the paper.
\begin{remark}[The extension to general $k$ is not coordinatewise]\label{rem:ld-soundness-simultaneity}
	The extension from $k = 1$ to general $k$ cited in the proof of Theorem~\ref{lem:ld-soundness} is the reduction deriving Theorem~4.43 from Theorem~4.40 in \cite{Natarajan2018NEEXP}: the $k$ answer polynomials are combined into the single polynomial $(\alpha, u) \mapsto \sum_i \alpha_i g_i(u)$ in $m + k$ variables, the case $k = 1$ is applied once to the combined strategy, and the $k$ polynomials are recovered from the resulting measurement by exact linearity in $\alpha$ and the Schwartz--Zippel lemma. It cannot be replaced by applying the case $k = 1$ to each coordinate separately and combining the coordinate measurements by the sandwich of Lemma~\ref{lem:ld-sandwich}: over $\F_{2^s}$ with $d = 1$ and $k = 2$ there are a joint point measurement and coordinate polynomial measurements, on the maximally entangled state, for which every coordinate relation of the theorem holds with error $0$ while every POVM with polynomial-pair outcomes has consistency defect at least $1 - 2/q$ against the joint point measurement; the computation is in \cite{gap:qpbt_ld-simultaneous-sandwich}. The formalization accordingly proves the polynomial-tuple conclusion of Theorem~\ref{lem:ld-soundness} for $k = 1$ from Theorem~\ref{thm:main-formal}, which is the only value instantiated in the proof of Lemma~\ref{lem:qld-4-7}; the general case awaits a formalization of the combining reduction.
	% The k = 1 conclusion is MIPStarRE.QPBT.exists_directSimultaneousPolynomialMeasurements_of_k_eq_one
	% (MIPStarRE/QPBT/Combining/DirectLowDegree/Transport/Simultaneous.lean); it is stated for the
	% directly indexed game of docs/paper-gaps/qpbt_ld-dimension-divisibility.tex, which has no
	% blueprint definition yet, so no \lean link is attached here.
\end{remark}

% The nodes of this subsection are formalization-support nodes recording the
% reduction cited in the proof of Theorem lem:ld-soundness; the paper states
% none of them separately.
\subsection{The combining reduction to simultaneity parameter one}
\label{sec:ld-combining-reduction}

The extension from $k = 1$ to general $k$ cited in the proof of
Theorem~\ref{lem:ld-soundness} is recorded here in the form in which it is
formalized.  The nodes below are formalization-support nodes: the source states
none of them separately, and together they constitute its derivation of
Theorem~4.43 from Theorem~4.40 in~\cite{Natarajan2018NEEXP}.  The game in
question is the directly indexed low-degree game of
\cite{gap:qpbt_ld-dimension-divisibility}, which samples the coordinate index of
a line directly instead of through a fiber of $\chi$; its point questions, line
questions, answer formats and win predicate are those of
Definition~\ref{def:ld-game}, and it is used here because the reduction applies
the case $k = 1$ at dimension $m + k$, where the divisibility condition of
Definition~\ref{def:ld-game} is not available.  A node carries $\leanok$ only
when the Lean declarations it names are complete and sorry-free.

\begin{definition}[Combined parameters of the reduction]
	\label{def:ld-combining-parameters}
	\lean{MIPStarRE.QPBT.DirectLdParams.combined}
	\leanok
	\uses{def:ld-game}
	For a parameter tuple $\mathsf{ldparams} = (q,m,d,k)$ as in
	Definition~\ref{def:ld-game}, the \emph{combined parameter tuple} is
	\[
		\mathsf{ldparams}^{\mathrm c} = (q,\; m+k,\; d,\; 1)\;:
	\]
	the dimension grows by the simultaneity parameter, the simultaneity
	parameter becomes $1$, and the field size, the degree and the admissibility
	of the field size are unchanged.  The degree is unchanged because the
	combining map of Definition~\ref{def:ld-combining-map} is linear in the
	fresh coordinates and therefore, in the individual-degree class of
	Definition~\ref{def:polyfunc}, raises no individual degree above $d$ as soon
	as $d \ge 1$.  The source works with total degree and passes from $d$ to
	$d+1$ at this step; that increment is an artifact of the total-degree
	formulation and does not occur here.
\end{definition}

\begin{definition}[The combining map on polynomial tuples]
	\label{def:ld-combining-map}
	\lean{MIPStarRE.QPBT.combinePolyTuple, MIPStarRE.QPBT.combinePolyTuple_eval}
	\lean{MIPStarRE.QPBT.combinePolyTuple_mem_polyFunc}
	\lean{MIPStarRE.QPBT.directCombinedPolynomial}
	\lean{MIPStarRE.QPBT.directCombinedPolynomial_eval}
	\leanok
	\uses{def:polyfunc, def:ld-meas, def:ld-combining-parameters}
	Coordinates of $\F_q^{m+k}$ are written $(u,\alpha)$, with $u \in \F_q^m$
	occupying the first $m$ coordinates and $\alpha \in \F_q^k$ the last $k$;
	the coordinates carrying $\alpha$ are called the \emph{combining
	coordinates}.  For a tuple $g = (g_0,\ldots,g_{k-1})$ of polynomials
	$g_r \in \polyfunc{m}{q}{d}$, the \emph{combined polynomial} is
	\[
		\mathrm{comb}_g(u,\alpha) \;=\; \sum_{r=0}^{k-1} \alpha_r\, g_r(u)\;.
	\]
	It lies in $\polyfunc{m+k}{q}{d}$ whenever $d \ge 1$: a component $g_r$
	involves no combining coordinate, so multiplying it by $\alpha_r$ leaves its
	individual degrees in the first $m$ coordinates unchanged and produces
	individual degree $1$ in the coordinate carrying $\alpha_r$.

	The order of the two coordinate blocks is part of the construction.  The
	point coordinates precede the combining coordinates, so that zeroing a
	prefix of a direction of $\F_q^{m+k}$, as the diagonal-line branch of
	Definition~\ref{def:ld-question-distribution} does, restricts on the first
	$m$ coordinates to zeroing the corresponding prefix of the point part of
	that direction; this is what makes the diagonal-line questions of the
	combined game project to diagonal-line questions of the original game in
	Definition~\ref{def:ld-combined-strategy}.  Definition~\ref{def:combine-map}
	is a variant of the same construction with two components in \emph{disjoint}
	blocks of point coordinates; the map above keeps one common block of point
	coordinates, as the reduction requires.
\end{definition}

\begin{lemma}[Recovery of the components of a combined polynomial]
	\label{lem:ld-combining-split}
	\lean{MIPStarRE.QPBT.splitCombinedPoly}
	\lean{MIPStarRE.QPBT.splitCombinedPoly_combinePolyTuple}
	\lean{MIPStarRE.QPBT.combinePolyTuple_injective}
	\lean{MIPStarRE.QPBT.IsDirectCombined, MIPStarRE.QPBT.directTupleOfCombined}
	\lean{MIPStarRE.QPBT.directTupleOfCombined_directCombinedPolynomial}
	\leanok
	\uses{def:ld-combining-map}
	Let $\mathrm{split}_r$ denote the substitution of the $r$-th standard basis
	vector of $\F_q^k$ for the combining coordinates, a ring homomorphism
	$\F_q[u,\alpha] \to \F_q[u]$ fixing the point coordinates.  Then
	$\mathrm{split}_r(\mathrm{comb}_g) = g_r$ for every tuple $g$ and every
	$r$.  Consequently $g \mapsto \mathrm{comb}_g$ is injective, and the map
	sending $p \in \polyfunc{m+k}{q}{d}$ to the unique tuple $g$ with
	$\mathrm{comb}_g = p$ when $p$ is of that form, and to the zero tuple
	otherwise, is a well-defined total map $\mathrm{split}$ with
	$\mathrm{split}(\mathrm{comb}_g) = g$.
\end{lemma}
\begin{proof}
	\leanok
	\uses{def:ld-combining-map}
	Applying the homomorphism $\mathrm{split}_r$ to the sum defining
	$\mathrm{comb}_g$ sends the summand of index $s$ to $g_s$ when $s = r$ and to
	$0$ otherwise, since the combining coordinate carrying $\alpha_s$ is sent to
	$1$ for $s = r$ and to $0$ otherwise, while $g_s$, which involves no
	combining coordinate, is fixed; the sum is therefore $g_r$.
	Injectivity follows, and it makes the case distinction defining
	$\mathrm{split}$ unambiguous.  The value on polynomials that are not
	combined is immaterial: the source leaves it arbitrary, and the estimates
	below use only the value on combined polynomials.
\end{proof}

\begin{lemma}[Collision of two linear forms in the combining variables]
	\label{lem:ld-combining-collision}
	\lean{MIPStarRE.QPBT.combiningLinearForm}
	\lean{MIPStarRE.QPBT.combiningLinearForm_injective}
	\lean{MIPStarRE.QPBT.combiningLinearForm_agreementProbability_le}
	\lean{MIPStarRE.QPBT.directCombiningCollision_avg_le}
	\leanok
	\uses{lem:schwartz-zippel-total-degree}
	Let $c, b \in \F_q^k$ with $c \neq b$.  Then
	\[
		\Pr_{\alpha \in \F_q^k}\Bigl[\; \sum_{r} c_r \alpha_r
			= \sum_{r} b_r \alpha_r \;\Bigr] \;\le\; \frac1q\;,
	\]
	the probability being over a uniformly random $\alpha$.
\end{lemma}
\begin{proof}
	\leanok
	\uses{lem:schwartz-zippel-total-degree}
	The two sides are the evaluations at $\alpha$ of the linear forms
	$\sum_r c_r X_r$ and $\sum_r b_r X_r$ in $k$ variables.  The coefficient of
	$X_r$ in such a form is $c_r$, so the two forms are distinct polynomials,
	and each has total degree at most $1$.  Lemma~\ref{lem:schwartz-zippel-total-degree}
	bounds the probability that two distinct polynomials of total degree at most
	$1$ agree at a uniformly random point of $\F_q^k$ by $1/q$.
\end{proof}

\begin{definition}[The combined strategy]
	\label{def:ld-combined-strategy}
	\lean{MIPStarRE.QPBT.directCombinedMeasuredQuestion}
	\lean{MIPStarRE.QPBT.directCombinedAnswerMap}
	\lean{MIPStarRE.QPBT.directCombinedStrategy}
	\lean{MIPStarRE.QPBT.directCombinedStrategy_isProjective}
	\leanok
	\uses{def:ld-game, def:ld-question-distribution, def:ld-win-predicate,
		def:projective-strategy-general, def:ld-combining-parameters,
		def:ld-combining-map}
	Let $\mathcal S = (\psi, A, B)$ be a projective strategy for the directly
	indexed low-degree game with parameters $\mathsf{ldparams}$.  The
	\emph{combined strategy} $\mathcal S^{\mathrm c}$ for the directly indexed
	low-degree game with parameters $\mathsf{ldparams}^{\mathrm c}$ has the same
	state and the same two Hilbert spaces, and answers each question of the
	combined game by measuring one question of $\mathcal S$ and post-processing
	its outcome.  Writing questions of the combined game in the coordinates of
	Definition~\ref{def:ld-combining-map}, the five cases are the following.
	\begin{enumerate}
		\item On the point question $(u,\alpha) \in \F_q^{m+k}$: measure the
		point question $u$, obtaining $b \in \F_q^k$, and answer
		$\sum_r \alpha_r b_r$.
		\item On the axis-parallel line question with stored index $i < m$ and
		canonical base $(u_0,\alpha_0)$: measure the axis-parallel line question
		with index $i$ and base $u_0$, obtaining univariate polynomials
		$f_0,\ldots,f_{k-1}$ of degree at most $d$, and answer
		$\sum_r (\alpha_0)_r f_r$, again of degree at most $d$.
		\item On the axis-parallel line question with stored index $m + r$ and
		canonical base $(u_0,\alpha_0)$, a line along which $u$ is constant and
		only the coordinate $\alpha_r$ varies: measure the point question $u_0$,
		obtaining $b$, and answer the degree-one polynomial
		$t \mapsto \sum_s (\alpha_0)_s b_s + t\, b_r$.
		\item On the diagonal line question with stored index $i < m$, canonical
		base $(u_0,\alpha_0)$ and direction $(v,w)$: measure the diagonal line
		question with index $i$, direction $v$ and the canonical base of the
		line $\ell(u_0, v)$, rebase the answer along the affine change of
		parameter relating the two presentations of that line, obtaining
		univariate polynomials $f_0,\ldots,f_{k-1}$ of degree at most $md$ in the
		parameter of the combined line, and answer
		$t \mapsto \sum_r \bigl((\alpha_0)_r + t\, w_r\bigr) f_r(t)$, of degree
		at most $md+1 \le (m+k)d$.
		\item On the diagonal line question with stored index $i \ge m$,
		canonical base $(u_0,\alpha_0)$ and direction $(0,w)$, a line along which
		$u$ is constant: measure the point question $u_0$, obtaining $b$, and
		answer $t \mapsto \sum_r (\alpha_0)_r b_r + t \sum_r w_r b_r$.
	\end{enumerate}
	Each case measures a single measurement of $\mathcal S$ and relabels its
	outcome, so $\mathcal S^{\mathrm c}$ is again a projective strategy, on the
	same state.  Cases~3 and~5 use that a line of $\F_q^{m+k}$ whose direction
	has vanishing point part is contained in a fiber of the projection to
	$\F_q^m$, so that the combined polynomial restricted to it is determined by
	the $k$ values of the tuple at the single point $u_0$; the rebasing in
	case~4 is the operation of \texttt{Transport.Questions}, needed because the
	canonical base of a line of $\F_q^{m+k}$ need not project to the canonical
	base of its image line in $\F_q^m$.
\end{definition}

\begin{lemma}[The questions measured by the combined strategy]
	\label{lem:ld-combined-question-law}
	\lean{MIPStarRE.QPBT.directCombined_measuredQuestion_of_pointVar}
	\lean{MIPStarRE.QPBT.directCombined_measuredQuestion_of_coefficientVar}
	\lean{MIPStarRE.QPBT.directCombinedSample_pointParts_uniform}
	\lean{MIPStarRE.QPBT.directCombinedSample_index_uniform}
	\lean{MIPStarRE.QPBT.directCombinedSample_indexPointParts_uniform}
	\leanok
	\uses{def:ld-question-distribution, def:ld-combined-strategy}
	Sample a question pair of the directly indexed game with parameters
	$\mathsf{ldparams}^{\mathrm c}$, and let $\mathbf{i}$ be the coordinate index
	of its common sample, uniform on $\{0,1,\ldots,m+k-1\}$.  Conditioned on
	$\mathbf{i} < m$, an event of probability $m/(m+k)$, the pair of questions
	measured by $\mathcal S^{\mathrm c}$ in Definition~\ref{def:ld-combined-strategy}
	is distributed as the question pair of the directly indexed game with
	parameters $\mathsf{ldparams}$ conditioned on the same ordered type pair.
	Conditioned on $\mathbf{i} \ge m$, both measured questions are the point
	question at the point part of the sampled point, and that point part is
	uniform on $\F_q^m$; that is, the measured pair is a point-versus-point pair
	of the game with parameters $\mathsf{ldparams}$.
\end{lemma}
\begin{proof}
	\leanok
	\uses{def:ld-question-distribution, def:ld-combining-map}
	The common sample of the combined game consists of a uniform point
	$(u,\alpha) \in \F_q^{m+k}$, a uniform index $\mathbf{i}$, and a uniform
	direction $(v,w)$, the three independent.  Conditioned on $\mathbf{i} < m$,
	the triple $(u, \mathbf{i}, v)$ is uniform on
	$\F_q^m \times \{0,\ldots,m-1\} \times \F_q^m$, which is the common sample of
	the game with parameters $\mathsf{ldparams}$, and $\alpha$ and $w$ are
	independent of it.  Under this conditioning the coordinate direction of an
	axis-parallel question and the prefix projection of a diagonal direction have
	point parts equal to the coordinate direction and the prefix projection
	computed in $\F_q^m$, by the coordinate order fixed in
	Definition~\ref{def:ld-combining-map}, and cases~2 and~4 of
	Definition~\ref{def:ld-combined-strategy} measure exactly the canonical
	question of the game with parameters $\mathsf{ldparams}$ attached to that
	sample.  Conditioned on $\mathbf{i} \ge m$, every direction produced by the
	question distribution has vanishing point part, so cases~1, 3 and~5 apply and
	all measure the point question at $u$, which is uniform.

	The two conditionings are legitimate because the coordinate index and
	the two point parts of the common sample are jointly uniform on
	$\{0,\ldots,m+k-1\} \times \F_q^m \times \F_q^m$, so the index is
	independent of the point parts.
\end{proof}

\begin{lemma}[Value of the combined strategy]
	\label{lem:ld-combined-value}
	\lean{MIPStarRE.QPBT.directLdRejectionProbability_directCombinedStrategy_le}
	\lean{MIPStarRE.QPBT.directCombinedStrategy_value_ge}
	\leanok
	\uses{def:ld-combined-strategy, lem:ld-combined-question-law,
		def:ld-win-predicate, def:tensor-product-value}
	There is a universal constant $C$ such that the following holds.  If
	$\mathcal S$ succeeds with probability at least $1 - \eps$ in the directly
	indexed low-degree game with parameters $\mathsf{ldparams}$, then
	$\mathcal S^{\mathrm c}$ succeeds with probability at least $1 - C\eps$ in the
	directly indexed low-degree game with parameters
	$\mathsf{ldparams}^{\mathrm c}$.
\end{lemma}
\begin{proof}
	\leanok
	\uses{def:ld-combined-strategy, lem:ld-combined-question-law,
		def:ld-win-predicate}
	Both games weight their nine ordered type pairs uniformly, so it suffices to
	bound the rejection probability of each branch of the combined game by a
	bounded number of branch rejection probabilities of the original game and to
	sum.  Fix a branch and split according to the event of
	Lemma~\ref{lem:ld-combined-question-law}.

	On the event $\mathbf{i} < m$ the two combined answers are the images of the
	two original answers under the relabelings of
	Definition~\ref{def:ld-combined-strategy}, and each clause of
	Definition~\ref{def:ld-win-predicate} for the combined game follows from the
	corresponding clause for the original game.  For the point-versus-point
	clause this is immediate.  For a line-versus-point clause, let $t$ present
	the sampled combined point on the sampled combined line; then $t$ presents
	the point part of the sampled point on the image line, the combining part of
	the sampled point is $\alpha_0 + t w$ in the diagonal case and $\alpha_0$ in
	the axis-parallel case, and the combined line answer evaluated at $t$ is
	$\sum_r (\alpha_0 + t w)_r f_r(t)$, respectively $\sum_r (\alpha_0)_r f_r(t)$;
	substituting the original clause $f_r(t) = b_r$ for every $r$ turns this into
	the combined point answer $\sum_r \alpha_r b_r$.  For a line-versus-line
	clause, equality of the two original tuples of line answers gives equality of
	their combinations.

	On the event $\mathbf{i} \ge m$ all measured questions are the point question
	at the same point $u$, and every clause of the combined game follows from the
	agreement $b = b'$ of the two original point answers, which is the
	point-versus-point clause of the original game: in cases~3 and~5 the combined
	line answer is a fixed affine function of the measured tuple, and evaluating
	it at the parameter presenting the sampled combined point returns
	$\sum_r \alpha_r b_r$, the combined point answer computed from the same
	tuple.

	Each branch of the combined game is therefore rejected with probability at
	most the sum of two branch rejection probabilities of the original game, and
	summing over the nine branches bounds the rejection probability of
	$\mathcal S^{\mathrm c}$ by an absolute multiple of that of $\mathcal S$.
\end{proof}

\begin{lemma}[Exact linearity in the combining variables]
	\label{lem:ld-combining-exact-linearity}
	\lean{MIPStarRE.QPBT.combinedRestrict}
	\lean{MIPStarRE.QPBT.eval_combinedRestrict}
	\lean{MIPStarRE.QPBT.combinedRestrict_combinePolyTuple}
	\lean{MIPStarRE.QPBT.combinedRestrict_eq_of_eq_combiningLinearForm}
	\lean{MIPStarRE.QPBT.combinedCoef}
	\lean{MIPStarRE.QPBT.combinePolyTuple_combinedCoef_iff}
	\lean{MIPStarRE.QPBT.exists_eq_combiningLinearForm_combinedRestrict_iff}
	\lean{MIPStarRE.QPBT.totalDegree_combinedCoef_le}
	\lean{MIPStarRE.QPBT.combinedRestrict_combiningLinearForm_avg_le}
	\lean{MIPStarRE.QPBT.directCombinedExactLinearity_avg_le}
	\leanok
	\uses{def:ld-combining-map, lem:ld-combining-split,
		lem:schwartz-zippel-total-degree}
	Let $p \in \polyfunc{m+k}{q}{d}$ be a polynomial which is not the combined
	polynomial of any tuple.  Then, for a uniformly random $u \in \F_q^m$, the
	polynomial $p(u,\cdot) \in \F_q[\alpha_0,\ldots,\alpha_{k-1}]$ obtained by
	substituting $u$ for the point coordinates is a linear form
	$\sum_r c_r \alpha_r$ with probability at most $md/q$.
\end{lemma}
\begin{proof}
	\leanok
	\uses{lem:schwartz-zippel-total-degree, def:ld-combining-map,
		lem:ld-combining-split}
	Write $p = \sum_{\mu} c_\mu \, \alpha^\mu$, the sum over exponent vectors
	$\mu \in \N^k$ of the combining coordinates, with coefficients
	$c_\mu \in \polyfunc{m}{q}{d}$.  By
	Definition~\ref{def:ld-combining-map} and
	Lemma~\ref{lem:ld-combining-split}, $p$ is a combined polynomial precisely
	when $c_\mu = 0$ for every $\mu$ outside the $k$ standard basis vectors, and
	$p(u,\cdot)$ is a linear form precisely when $c_\mu(u) = 0$ for every such
	$\mu$.  If $p$ is not combined, fix one such $\mu$ with $c_\mu \neq 0$; the
	total degree of $c_\mu$ is at most $md$, so
	Lemma~\ref{lem:schwartz-zippel-total-degree} bounds the probability of
	$c_\mu(u) = 0$ by $md/q$.
\end{proof}

\begin{lemma}[Recovery of the polynomial-tuple measurements]
	\label{lem:ld-combining-recovery}
	\uses{def:ld-meas, def:consistency, def:bracket, def:ld-combining-map,
		lem:ld-combining-split, lem:ld-combining-collision,
		lem:ld-combining-exact-linearity, def:ld-combined-strategy}
	Let $\mathcal S$ be a projective strategy for the directly indexed game with
	parameters $\mathsf{ldparams}$, and let $G^{\mathrm A}$, $G^{\mathrm B}$ be
	measurements with outcomes in $\polyfunc{m+k}{q}{d}$ satisfying, for the
	combined strategy $\mathcal S^{\mathrm c}$ and with error $\delta$, the three
	conclusions of Theorem~\ref{lem:ld-soundness} at parameters
	$\mathsf{ldparams}^{\mathrm c}$.  Let $H^w$ be $G^w$ post-processed by the
	map $\mathrm{split}$ of Lemma~\ref{lem:ld-combining-split}, so that
	$H^w \in \mathrm{PolyMeas}(m,q,d,k)$.  Then the three conclusions of
	Theorem~\ref{lem:ld-soundness} hold for $\mathcal S$ at parameters
	$\mathsf{ldparams}$, with the measurements $H^{\mathrm A}$, $H^{\mathrm B}$
	and with error $O(\delta + (m+k)d/q)$.
\end{lemma}
\begin{proof}
	\uses{lem:ld-combining-collision, lem:ld-combining-exact-linearity,
		lem:ld-combining-split, def:bracket, lem:schwartz-zippel-total-degree}
	The first conclusion for $\mathcal S^{\mathrm c}$ says that, on average over
	a uniform $(u,\alpha) \in \F_q^{m+k}$, the value $p(u,\alpha)$ of the outcome
	$p$ of $G^{\mathrm B}$ agrees with the combined point answer
	$\sum_r \alpha_r b_r$ computed from the outcome $b$ of the point measurement
	of $\mathcal S$ at $u$, up to $\delta$.

	The outcome is a combined polynomial with probability at least
	$1 - O(\delta + (m+k)d/q)$.  Indeed, condition on an outcome $p$ that is not
	combined.  By Lemma~\ref{lem:ld-combining-exact-linearity}, $p(u,\cdot)$ fails
	to be a linear form except with probability $md/q$ over $u$; the combined
	point answer is a linear form in $\alpha$ by construction; and two distinct
	polynomials in $\alpha$ of total degree at most $kd$ agree at a uniform
	$\alpha$ with probability at most $kd/q$ by
	Lemma~\ref{lem:schwartz-zippel-total-degree}.  Hence, conditioned on a
	non-combined outcome, the agreement above holds with probability at most
	$(m+k)d/q$, and comparing with the bound $\delta$ gives the claim.

	Replacing $G^w$ by $H^w$ therefore changes each of the first two relations by
	at most $O(\delta + (m+k)d/q)$, since the two measurements differ only on the
	event that the outcome is not combined, and
	$\mathrm{split}(\mathrm{comb}_g) = g$ on its complement; the third relation
	for $H^{\mathrm A}$, $H^{\mathrm B}$ follows from the third relation for
	$G^{\mathrm A}$, $G^{\mathrm B}$ by data processing, $H^w$ being a
	post-processing of $G^w$.

	It remains to upgrade the first relation from the combined value to the tuple
	of values.  Let $g$ be the outcome of $H^{\mathrm B}$ and $b$ the outcome of
	the point measurement at $u$.  If $(g_0(u),\ldots,g_{k-1}(u)) \neq b$, then by
	Lemma~\ref{lem:ld-combining-collision} the combined values
	$\sum_r \alpha_r g_r(u)$ and $\sum_r \alpha_r b_r$ agree with probability at
	most $1/q$ over a uniform $\alpha$.  Averaging over $\alpha$ therefore
	converts the relation for the combined values into the tuple relation
	\[
		A^{\mathsf{Point},\,u}_{a_1,\ldots,a_k} \ot \Id_{\mathrm B}
		\simeq \Id_{\mathrm A} \ot
			H^{\mathrm B}_{[(g_1(u),\ldots,g_k(u)) = (a_1,\ldots,a_k)]}
	\]
	with error $O(\delta + 1/q)$, and symmetrically for the second relation.
\end{proof}

\begin{proposition}[Simultaneous polynomial measurements for general $k$]
	\label{prop:ld-simultaneous-general-k}
	\uses{def:ld-meas, def:consistency, def:bracket,
		def:projective-strategy-general, def:ld-combining-parameters,
		lem:ld-combined-value, lem:ld-combining-recovery}
	Let $\mathsf{ldparams} = (q,m,d,k)$ and let $\mathcal S = (\psi, A, B)$ be a
	projective strategy succeeding with probability at least $1 - \eps$ in the
	directly indexed low-degree game with parameters $\mathsf{ldparams}$.  Then
	there exist measurements $G^{\mathrm A}, G^{\mathrm B} \in
	\mathrm{PolyMeas}(m,q,d,k)$ satisfying the three consistency relations of
	Theorem~\ref{lem:ld-soundness} with an error of the form
	$\delta_{\mathrm{ld}}(\eps, q, m+k, d, 1) + O((m+k)d/q)$, which is absorbed
	into $\delta_{\mathrm{ld}}(\eps,q,m,d,k)$ by enlarging its universal
	constants.
\end{proposition}
\begin{proof}
	\uses{lem:ld-combined-value, lem:ld-combining-recovery,
		def:ld-combining-parameters}
	Apply Lemma~\ref{lem:ld-combined-value} to obtain a projective strategy
	$\mathcal S^{\mathrm c}$ for the directly indexed game with parameters
	$\mathsf{ldparams}^{\mathrm c}$ succeeding with probability at least
	$1 - C\eps$.  That game has simultaneity parameter $1$, so the case $k = 1$
	of Theorem~\ref{lem:ld-soundness} applies and produces measurements
	$G^{\mathrm A}$, $G^{\mathrm B}$ with outcomes in $\polyfunc{m+k}{q}{d}$
	satisfying the three relations with error
	$\delta_{\mathrm{ld}}(C\eps, q, m+k, d, 1)$.  Lemma~\ref{lem:ld-combining-recovery}
	converts them into measurements in $\mathrm{PolyMeas}(m,q,d,k)$ satisfying
	the three relations for $\mathcal S$ with the stated error.
\end{proof}

\section{The Magic Square game}
% Paper label sec:ms kept for cross-reference fidelity.
\label{sec:ms}

We recall the Magic Square game of Mermin and Peres. It is a simple self-test for two EPR pairs, and it allows one to test that a pair of binary observables \emph{anticommutes}; both properties are used in the Pauli basis test. Here it is presented as a binary constraint satisfaction game. Throughout this section and the next, $\sigma^X$ and $\sigma^Z$ denote the single-qubit Pauli observables, with associated two-outcome projective measurements $\{\sigma^W_b\}_{b \in \F_2}$ for $W \in \{X,Z\}$.

\begin{definition}[The Magic Square game]\label{def:ms-game}
	\uses{def:game, def:graph-distribution}
	The Magic Square game $\mathfrak{G}^{\mathrm{MS}}$ is the following two-player game. There are $9$ variables $x_1,\ldots,x_9$ taking values in $\F_2$, arranged in a $3 \times 3$ grid: the cell in row $r$ and column $c$ contains the variable $x_{3(r-1)+c}$. There are $6$ linear equations over these variables: for each of the three rows and each of the first two columns, the sum of the three variables in it equals $0$; for the third column, the sum of the three variables in it equals $1$.

	The question types are
	\begin{align*}
		\mathcal T^{\mathrm{mc}} &= \{\mathsf{Constraint}_i : i = 1,2,\ldots,6\}\;,\\
		\mathcal T^{\mathrm{mv}} &= \{\mathsf{Variable}_j : j = 1,2,\ldots,9\}\;,\\
		\mathcal T^{\mathrm{ms}} &= \mathcal T^{\mathrm{mc}} \cup \mathcal T^{\mathrm{mv}}\;,
	\end{align*}
	where $\mathsf{Constraint}_i$ for $i \in \{1,2,3\}$ denotes the constraint of row $i$ (with variables $x_{3(i-1)+1}, x_{3(i-1)+2}, x_{3(i-1)+3}$), $\mathsf{Constraint}_4$ and $\mathsf{Constraint}_5$ denote the constraints of the first two columns, and $\mathsf{Constraint}_6$ denotes the constraint of the third column (with variables $x_{i-3}, x_{i}, x_{i+3}$ for $\mathsf{Constraint}_i$, $i \in \{4,5,6\}$).

	The question distribution is as follows: the verifier samples $\mathsf{Constraint}_i$ uniformly from $\mathcal T^{\mathrm{mc}}$, then samples $\mathsf{Variable}_j$ uniformly among the three variables appearing in the row or column of $\mathsf{Constraint}_i$; one player, chosen uniformly at random, receives $\mathsf{Constraint}_i$ and the other receives $\mathsf{Variable}_j$. Equivalently, the verifier samples a uniformly random edge of the bipartite \emph{type graph} $G^{\mathrm{ms}}$ on vertex sets $\mathcal T^{\mathrm{mc}}$ and $\mathcal T^{\mathrm{mv}}$, whose edges are $\mathsf{Constraint}_i \,\text{---}\, \mathsf{Variable}_{3(i-1)+j}$ for $i,j \in \{1,2,3\}$ and $\mathsf{Constraint}_i \,\text{---}\, \mathsf{Variable}_{(i-3)+3(j-1)}$ for $i \in \{4,5,6\}$, $j \in \{1,2,3\}$ ($18$ edges in total, and no self-loops), and sends one endpoint to a uniformly random player and the other endpoint to the other player.

	The $\mathsf{Constraint}$ player answers with three bits $(\beta_{v_1}, \beta_{v_2}, \beta_{v_3}) \in \F_2^3$, where $(v_1,v_2,v_3)$ are the indices of the three variables of $\mathsf{Constraint}_i$; the $\mathsf{Variable}$ player answers with a single bit $\gamma \in \F_2$. The players win if and only if the $\mathsf{Constraint}$ player's answer satisfies the equation associated with $\mathsf{Constraint}_i$ and $\gamma = \beta_j$.
\end{definition}

The following theorem records the self-testing (rigidity) properties of the Magic Square game. It is used in the soundness analysis of the Pauli basis test to enforce anticommutation relations between certain pairs of observables.
With zero-based labels, the paper's variables $1$ and $5$ have labels $0$ and
$4$, respectively.

\begin{theorem}[Rigidity of the Magic Square game]\label{thm:ms-rigidity}
	\lean{MIPStarRE.QPBT.exists_ms_rigidity}
	\leanok
	\uses{def:ms-game, def:tensor-product-value, def:EPR, def:povm-distance}
	Let $\mathcal S = (\ket{\psi}, A, B)$ be a strategy that succeeds in the Magic Square game $\mathfrak{G}^{\mathrm{MS}}$ with probability $1-\eps$, where $\ket{\psi} \in \hilb_{\mathrm A} \ot \hilb_{\mathrm B}$. Then there exist local isometries $\phi_{\mathrm A}\colon \hilb_{\mathrm A} \to \hilb_{\mathrm A'} \ot \hilb_{\mathrm A''}$ and $\phi_{\mathrm B}\colon \hilb_{\mathrm B} \to \hilb_{\mathrm B'} \ot \hilb_{\mathrm B''}$, where $\hilb_{\mathrm A'}, \hilb_{\mathrm B'} \cong (\C^2)^{\ot 2}$ and $\hilb_{\mathrm A''}, \hilb_{\mathrm B''}$ are finite dimensional, and a state $\ket{\mathrm{aux}} \in \hilb_{\mathrm A''} \ot \hilb_{\mathrm B''}$, such that:
	\begin{enumerate}
		\item $\bigl\| \phi_{\mathrm A} \ot \phi_{\mathrm B} \ket{\psi} - \ket{\mathrm{EPR}_2}^{\ot 2} \ot \ket{\mathrm{aux}} \bigr\| \le O(\sqrt\eps)$;
		\item letting $\tilde A^x_a = \phi_{\mathrm A} A^x_a \phi_{\mathrm A}^\dagger$ and $\tilde B^y_b = \phi_{\mathrm B} B^y_b \phi_{\mathrm B}^\dagger$,
		\begin{align*}
			\tilde A^{\mathsf{Variable}_1}_b \ot I_{\mathrm{B'B''}} &\approx_{\sqrt\eps} (\sigma^X_b)_{\mathrm A'} \ot I_{\mathrm{A''B'B''}}\;,\\
			\tilde A^{\mathsf{Variable}_5}_b \ot I_{\mathrm{B'B''}} &\approx_{\sqrt\eps} (\sigma^Z_b)_{\mathrm A'} \ot I_{\mathrm{A''B'B''}}\;,\\
			I_{\mathrm{A'A''}} \ot \tilde B^{\mathsf{Variable}_1}_b &\approx_{\sqrt\eps} I_{\mathrm{A'A''B''}} \ot (\sigma^X_b)_{\mathrm B'}\;,\\
			I_{\mathrm{A'A''}} \ot \tilde B^{\mathsf{Variable}_5}_b &\approx_{\sqrt\eps} I_{\mathrm{A'A''B''}} \ot (\sigma^Z_b)_{\mathrm B'}\;,
		\end{align*}
		where each $\approx_{\sqrt\eps}$ holds with respect to the state $\ket{\mathrm{EPR}_2}^{\ot 2}_{\mathrm{A'B'}} \ot \ket{\mathrm{aux}}_{\mathrm{A''B''}}$ and the answer summation is over $b \in \{0,1\}$, and the approximation bounds carry implicit universal constants.
	\end{enumerate}
	As a consequence, letting $\tilde A^{\mathsf{Variable}_1} = \tilde A^{\mathsf{Variable}_1}_0 - \tilde A^{\mathsf{Variable}_1}_1$ and $\tilde A^{\mathsf{Variable}_5} = \tilde A^{\mathsf{Variable}_5}_0 - \tilde A^{\mathsf{Variable}_5}_1$ denote the associated observables, it holds that
	\begin{align*}
		\tilde A^{\mathsf{Variable}_1} \tilde A^{\mathsf{Variable}_5} \ot I_{\mathrm{B'B''}} &\approx_{\sqrt\eps} -\, \tilde A^{\mathsf{Variable}_5} \tilde A^{\mathsf{Variable}_1} \ot I_{\mathrm{B'B''}}\;,\\
		I_{\mathrm{A'A''}} \ot \tilde B^{\mathsf{Variable}_1} \tilde B^{\mathsf{Variable}_5} &\approx_{\sqrt\eps} -\, I_{\mathrm{A'A''}} \ot \tilde B^{\mathsf{Variable}_5} \tilde B^{\mathsf{Variable}_1}\;,
	\end{align*}
	where, for operators $M$, $N$ that are not POVM elements, the notation $M \approx_\delta N$ means $\bra{\psi} (M-N)^\dagger (M-N) \ket{\psi} \le O(\delta)$ with respect to the same state.
\end{theorem}
\begin{proof}
	This result is imported from Theorem~6.9 of~\cite{Coladangelo2017Robust}. There are two minor differences between their statement and the one given here. First, they state a robustness of $O(\eps)$, but measure the closeness of the actual and ideal states in trace norm; expressed as the Euclidean distance between $\phi_{\mathrm A} \ot \phi_{\mathrm B}\ket{\psi}$ and $\ket{\mathrm{EPR}_2}^{\ot 2} \ot \ket{\mathrm{aux}}$, this becomes $O(\sqrt\eps)$. Second, their ideal observables for the $\mathsf{Variable}_1$ and $\mathsf{Variable}_5$ questions are $I \ot \sigma^Z$ and $\sigma^Z\sigma^X \ot \sigma^Z\sigma^X$; under a local change of basis these are equivalent to $\sigma^X \ot I$ and $\sigma^Z \ot I$, respectively.
\end{proof}

The converse construction shows that any pair of anticommuting binary observables yields a perfect strategy for the Magic Square game.

\begin{theorem}\label{thm:ms-from-ac}
	\lean{MIPStarRE.QPBT.obsOf, MIPStarRE.QPBT.exists_ms_perfect_strategy_of_anticommuting}
	\leanok
	\uses{def:ms-game, def:consistent-measurement, def:spcc, def:EPR, def:tensor-product-value}
	Let $A = \{A_b\}_{b \in \F_2}$ and $B = \{B_b\}_{b \in \F_2}$ be two-outcome projective measurements acting on $(\C^q)^{\ot n}$ that are consistent on $\ket{\mathrm{EPR}_q}^{\ot n}$ (Definition~\ref{def:consistent-measurement}), and let $\mathcal O_A = A_0 - A_1$ and $\mathcal O_B = B_0 - B_1$ be the corresponding $\pm 1$-valued observables. Suppose that $\mathcal O_A \mathcal O_B = - \mathcal O_B \mathcal O_A$. Then there exists a symmetric strategy $\mathcal S = (\psi, M)$ for the Magic Square game with the following properties:
	\begin{enumerate}
		\item $\mathcal S$ is an SPCC (symmetric, projective, consistent, commuting) strategy of value $1$;
		\item the state is $\ket{\psi} = \ket{\mathrm{EPR}_q}^{\ot n} \ot \ket{\mathrm{EPR}_2}$;
		\item for $b \in \{0,1\}$, $M^{\mathsf{Variable}_1}_b = A_b \ot I$ and $M^{\mathsf{Variable}_5}_b = B_b \ot I$.
	\end{enumerate}
\end{theorem}
\begin{proof}
	\uses{def:consistent-measurement}
	The strategy is based on the canonical two-qubit operator solution of the Magic Square, with the state $\ket{\psi} = \ket{\mathrm{EPR}_q}^{\ot n} \ot \ket{\mathrm{EPR}_2}$ and the cells of the grid assigned the $\pm 1$-valued observables
	\[
		\begin{array}{ccc}
			\mathcal O_A \ot I & I \ot \sigma^X & \mathcal O_A \ot \sigma^X \\[2pt]
			I \ot \sigma^Z & \mathcal O_B \ot I & \mathcal O_B \ot \sigma^Z \\[2pt]
			\mathcal O_A \ot \sigma^Z & \mathcal O_B \ot \sigma^X & \mathcal O_A \mathcal O_B \ot \sigma^Z \sigma^X
		\end{array}
	\]
	on $(\C^q)^{\ot n} \ot \C^2$: on question $\mathsf{Variable}_j$ a player measures the two-outcome projective measurement of the observable in cell $j$, and on question $\mathsf{Constraint}_i$ a player simultaneously performs the three measurements of the corresponding row or column; item~3 holds by construction. Anticommutation of $\mathcal O_A$ and $\mathcal O_B$ gives
	\begin{equation}\label{eq:am-i-hermitian}
		\mathcal O_A \mathcal O_B \ot \sigma^Z \sigma^X
		= - \mathcal O_A \mathcal O_B \ot \sigma^X \sigma^Z
		= \mathcal O_B \mathcal O_A \ot \sigma^X \sigma^Z
		= (\mathcal O_A \mathcal O_B \ot \sigma^Z \sigma^X)^\dagger\;,
	\end{equation}
	so the bottom-right cell is Hermitian, and it squares to the identity; hence every cell is a genuine $\pm 1$-valued observable. The three observables in each row and each column commute pairwise --- for the third row and third column this again uses Equation~\eqref{eq:am-i-hermitian} --- so the $\mathsf{Constraint}$ measurements are well defined, and the strategy is symmetric, projective, and commuting. A two-outcome projective measurement is consistent on a state if and only if its observable $\mathcal O$ satisfies $\mathcal O \ot I \ket{\psi} = I \ot \mathcal O \ket{\psi}$; since $A$ and $B$ are consistent on $\ket{\mathrm{EPR}_q}^{\ot n}$ by hypothesis and $\sigma^X$, $\sigma^Z$ are consistent on $\ket{\mathrm{EPR}_2}$, each cell observable is a product of consistent, commuting factors and is therefore consistent on $\ket{\psi}$, and consistency of the $\mathsf{Constraint}$ measurements follows from the cellwise consistencies together with commutativity. Finally, the observables in each row and each of the first two columns multiply to $I$, while those of the third column multiply to $-I$; for the measurement $\{S_0, S_1\}$ obtained by measuring a $\mathsf{Constraint}_i$ triple and outputting the sum of the three bits, the observable $S_0 - S_1$ equals the product of the three cell observables, hence $\pm I$, which forces the three answer bits to sum to $0$ (respectively, $1$ for $\mathsf{Constraint}_6$) with certainty. Together with consistency this shows that every check of the Magic Square game is passed with probability $1$, so the strategy has value $1$.
\end{proof}

\section{The Pauli basis test}
% Paper label sec:qld-game kept for cross-reference fidelity.
\label{sec:qld-game}

The Pauli basis test is the quantum low individual degree test of Natarajan and Vidick, in the slightly modified form introduced by Natarajan and Wright~\cite{Natarajan2018NEEXP}. Informally, the test certifies that the players share the state $\ket{\mathrm{EPR}_q}^{\ot M}$, where $M = 2^m$. On a question of type $(\mathsf{Pauli}, W)$ with $W \in \{X,Z\}$, a player is expected to measure the generalized Pauli basis POVM $\{\tau^W_h\}_{h \in \F_q^M}$ of Chapter~\ref{chap:qpbt-algebra} and report the outcome $h$. Most subtests probe limited information about $h$: given a point $u_W \in \F_q^m$, a player is expected to report the evaluation $g_h(u_W)$ of the low-degree encoding $g_h$ of $h$, and for a fixed basis $W$ these probes are checked against each other by the classical low individual degree game of Section~\ref{sec:ld-verifier}. Consistency \emph{between} the two bases is checked through two-outcome probes: given additionally $r_W \in \F_q$, a player reports $\tr(g_h(u_W)\, r_W) \in \F_2$. The two binary measurements obtained in this way for $W = X$ and $W = Z$ either commute or anticommute, according to the value $\gamma \in \F_2$ of Equation~\eqref{eq:gamma-value} below. When $\gamma = 0$ the verifier checks commutation directly by asking for both bits simultaneously; when $\gamma = 1$ the verifier checks anticommutation by playing the Magic Square game, using Theorem~\ref{thm:ms-from-ac}.

\begin{definition}[Admissible parameters]\label{def:admissible}
	\uses{def:admissible-size}
	% The source introduces q, m, d in \N with \N the positive integers
	% (references/qpbt-paper/04_preliminaries.tex:6 and
	% references/qpbt-paper/08_classical_and_quantum_low_degree_tests.tex:903-905),
	% so d >= 1 is implicit in its notion of admissibility (lines 958-961 there);
	% we state it explicitly. It is load-bearing: at d = 0 the soundness function
	% delta_qld of thm:pauli vanishes identically while every strategy satisfies
	% the success hypothesis at eps = 1, and the low-degree check of
	% def:pauli-win-predicate invokes the classical game at ldparams = (q,m,d,1),
	% whose domain (Definition def:ld-game) requires d >= 1. The hypothesis
	% propagates to lem:pauli-completeness, thm:pauli, and cor:pauli-binary
	% through their quantification over admissible tuples.
	The tuple $\mathsf{qldparams} = (q,m,d)$ is \emph{admissible} if $d \ge 1$, $q$ is an admissible field size (Definition~\ref{def:admissible-size}), and $m$ divides $q$. Since an admissible field size satisfies $q \ge 2$, the divisibility forces $m \ge 1$, so an admissible tuple has all entries positive.
\end{definition}

\begin{definition}[Question distribution of the Pauli basis test]\label{def:pauli-question-distribution}
	\uses{def:admissible, def:ld-question-distribution, def:ms-game, def:cl-func, def:cl-dist, def:graph-distribution}
	% Admissibility of the full tuple is required here, not only the divisibility m | q:
	% the source parametrizes the Pauli basis test by the tuple qldparams introduced
	% together with Definition def:admissible (references/qpbt-paper/
	% 08_classical_and_quantum_low_degree_tests.tex:958-1075), and both the fixed binary
	% identification of F_q used by the map chi and the trace tr in eq:gamma-value
	% require q to be a power of 2.
	The game $\mathfrak{G}^{\mathrm{Pauli}}$ (also $\mathfrak{G}^{\mathrm{Pauli}}_{\mathsf{qldparams}}$) is parametrized by an admissible tuple $\mathsf{qldparams} = (q,m,d)$ (Definition~\ref{def:admissible}): admissibility makes $q$ a power of $2$, so that the field $\F_q$, its fixed identification with binary strings (Chapter~\ref{chap:qpbt-algebra}), and the trace $\tr\colon \F_q \to \F_2$ used by the win predicate below are defined, and it provides the divisibility $m \mid q$ required by the index map $\chi$ of Definition~\ref{def:ld-question-distribution}. The question type set is
	\begin{equation}
		\label{eq:pauli-type}
		\mathcal T^{\mathrm{pauli}} = \bigl( \{\mathsf{Point}, \mathsf{ALine}, \mathsf{DLine}, \mathsf{Pauli}, \mathsf{Pair}\} \times \{X,Z\} \bigr) \cup \mathcal T^{\mathrm{ms}} \cup \{\mathsf{Pair}\}\;,
	\end{equation}
	where $\mathcal T^{\mathrm{ms}}$ is the type set of the Magic Square game (Definition~\ref{def:ms-game}).

	The ambient space is $V^{\mathrm{pauli}} = (\F_q^m)^2 \times \F_q \times \F_q^m \times (\F_q)^2$, decomposed into the register subspaces $V_X, V_Z$ (each $m$-dimensional over $\F_q$), $V_I$ ($1$-dimensional), $V_V$ ($m$-dimensional), and $V_{R_X}, V_{R_Z}$ ($1$-dimensional each); elements of $V^{\mathrm{pauli}}$ are written $(u_X, u_Z, s, v, r_X, r_Z)$. For each type $t \in \mathcal T^{\mathrm{pauli}}$ define a CL function $L_t\colon V^{\mathrm{pauli}} \to V^{\mathrm{pauli}}$ as follows, where $L_{\mathsf{Point}}, L_{\mathsf{ALine}}, L_{\mathsf{DLine}}$ are the CL functions of Definition~\ref{def:ld-question-distribution}:
	\begin{enumerate}
		\item for $W \in \{X,Z\}$, $L_{(\mathsf{Point},W)}(u_X,u_Z,s,v,r_X,r_Z) = L_{\mathsf{Point}}(u_W,s,v)$, a $1$-level CL function, where the range $V_W \oplus V_I \oplus V_V$ of $L_{\mathsf{Point}}$ is embedded in $V^{\mathrm{pauli}}$ in the natural way (the same convention applies in items~2 and~3);
		\item for $W \in \{X,Z\}$, $L_{(\mathsf{ALine},W)}(u_X,u_Z,s,v,r_X,r_Z) = L_{\mathsf{ALine}}(u_W,s,v)$, a $2$-level CL function;
		\item for $W \in \{X,Z\}$, $L_{(\mathsf{DLine},W)}(u_X,u_Z,s,v,r_X,r_Z) = L_{\mathsf{DLine}}(u_W,s,v)$, a $3$-level CL function;
		\item for $t \in \mathcal T^{\mathrm{ms}} \cup \{\mathsf{Pair}, (\mathsf{Pair},X), (\mathsf{Pair},Z)\}$, $L_t(u_X,u_Z,s,v,r_X,r_Z) = (u_X,u_Z,0,0,r_X,r_Z)$, the $1$-level CL function projecting onto $V_X \oplus V_Z \oplus V_{R_X} \oplus V_{R_Z}$;
		\item for $W \in \{X,Z\}$, $L_{(\mathsf{Pauli},W)} = 0$, the identically zero ($0$-level) CL function.
	\end{enumerate}
	The maps $L_{(\mathsf{Point},W)}$, $L_{(\mathsf{ALine},W)}$, $L_{(\mathsf{DLine},W)}$ act as the zero function on $V_{\overline W} \oplus V_{R_X} \oplus V_{R_Z}$, where $\overline W$ denotes the basis other than $W$.

	The \emph{type graph} $G^{\mathrm{pauli}}$ has vertex set $\mathcal T^{\mathrm{pauli}}$ and the following edges: all $18$ edges of the Magic Square type graph $G^{\mathrm{ms}}$ (Definition~\ref{def:ms-game}); the edges
	\begin{align*}
		&(\mathsf{ALine},X) \,\text{---}\, (\mathsf{Point},X)\;, &&(\mathsf{DLine},X) \,\text{---}\, (\mathsf{Point},X)\;, &&(\mathsf{Point},X) \,\text{---}\, (\mathsf{Pauli},X)\;,\\
		&(\mathsf{ALine},Z) \,\text{---}\, (\mathsf{Point},Z)\;, &&(\mathsf{DLine},Z) \,\text{---}\, (\mathsf{Point},Z)\;, &&(\mathsf{Point},Z) \,\text{---}\, (\mathsf{Pauli},Z)\;,\\
		&(\mathsf{Point},X) \,\text{---}\, \mathsf{Variable}_1\;, &&(\mathsf{Point},Z) \,\text{---}\, \mathsf{Variable}_5\;, &&\\
		&(\mathsf{Point},X) \,\text{---}\, (\mathsf{Pair},X)\;, &&(\mathsf{Pair},X) \,\text{---}\, \mathsf{Pair}\;, &&\\
		&(\mathsf{Point},Z) \,\text{---}\, (\mathsf{Pair},Z)\;, &&(\mathsf{Pair},Z) \,\text{---}\, \mathsf{Pair}\;; &&
	\end{align*}
	and, in addition, a self-loop at every vertex.

	% The type pair is drawn from the graph distribution of Definition
	% def:graph-distribution, matching the source's formal definition
	% (references/qpbt-paper/07_types.tex:84-93 together with
	% references/qpbt-paper/08_classical_and_quantum_low_degree_tests.tex:1115-1120).
	% Sampling an undirected edge uniformly and then ordering its
	% endpoints uniformly would be a different distribution here: it gives each self-loop
	% twice the relative mass of an oriented non-loop pair, and G^pauli has a self-loop
	% at every vertex.
	The question distribution of $\mathfrak{G}^{\mathrm{Pauli}}$ is the typed CL distribution on $\mathcal T^{\mathrm{pauli}} \times V^{\mathrm{pauli}} \times \mathcal T^{\mathrm{pauli}} \times V^{\mathrm{pauli}}$ obtained by sampling the type pair $(t_{\mathrm A}, t_{\mathrm B})$ from the graph distribution $\mu_{G^{\mathrm{pauli}}}$ of Definition~\ref{def:graph-distribution}, that is, uniformly at random among all ordered pairs of types whose underlying unordered pair $\{t_{\mathrm A}, t_{\mathrm B}\}$ is an edge of $G^{\mathrm{pauli}}$, self-loops included; sampling $z \in V^{\mathrm{pauli}}$ uniformly at random; and setting $x_w = L_{t_w}(z)$ for $w \in \{\mathrm A,\mathrm B\}$. Explicitly, if $z = (u_X,u_Z,s,v,r_X,r_Z)$, then a player whose type is $(\mathsf{Point},W)$ receives $u_W$; a player whose type is $(\mathsf{ALine},W)$ receives $(u_0, s)$ with $u_0 = L^{\mathrm{Ln}}_{e_i}(u_W)$ and $i = \chi(s)$; a player whose type is $(\mathsf{DLine},W)$ receives $(u_0, s, v')$ with $i = \chi(s)$, $v' = \pi_{i-1}(v)$, and $u_0 = L^{\mathrm{Ln}}_{v'}(u_W)$; a player whose type lies in $\mathcal T^{\mathrm{ms}} \cup \{\mathsf{Pair}, (\mathsf{Pair},X), (\mathsf{Pair},Z)\}$ receives $(u_X, u_Z, r_X, r_Z)$; and a player whose type is $(\mathsf{Pauli},W)$ receives $0$.
\end{definition}

\begin{lemma}[Conditionally linear Pauli question maps]\label{lem:pauli-question-cl-components}
	\lean{MIPStarRE.QPBT.isCondLinear_pauliCL}
	\uses{def:pauli-question-distribution, def:cl-func}
	For every Pauli question type $t \in \mathcal T^{\mathrm{pauli}}$, the map $L_t$ of Definition~\ref{def:pauli-question-distribution} is conditionally linear. Its levels are $1$ for the point maps, $2$ for the axis-parallel line maps, $3$ for the diagonal line maps, $1$ for the shared projection maps, and $0$ for the Pauli maps.
\end{lemma}
\begin{proof}
	\uses{lem:ld-point-level, lem:ld-aline-level, lem:ld-dline-level}
	Transport the point, axis-parallel line, and diagonal line representations to the selected basis register. The shared projection onto the point and auxiliary registers is linear, while the Pauli maps are zero. These representations have levels $1$, $2$, $3$, $1$, and $0$, respectively.
\end{proof}

\begin{lemma}[Common level of the Pauli question maps]\label{lem:pauli-question-cl-family}
	\lean{MIPStarRE.QPBT.isTypedCondLinearFamily_pauliCL}
	\uses{def:pauli-question-distribution, def:typed-cl-functions}
	For every Pauli question type $t \in \mathcal T^{\mathrm{pauli}}$, the map $L_t$ of Definition~\ref{def:pauli-question-distribution} is a $3$-level CL function. Thus $L = \{L_t\}_{t \in \mathcal T^{\mathrm{pauli}}}$ is a finite typed family of $3$-level CL functions.
\end{lemma}
\begin{proof}
	\uses{lem:pauli-question-cl-components, lem:cl-level-mono}
	The component maps have levels at most $3$ by Lemma~\ref{lem:pauli-question-cl-components}. Lemma~\ref{lem:cl-level-mono} raises each representation to the common level $3$.
\end{proof}

\begin{lemma}[Pauli question-distribution identity]\label{lem:pauli-question-typed-equality}
	\lean{MIPStarRE.QPBT.pauliQuestionDistribution_eq_typedCL}
	\uses{def:pauli-question-distribution, def:typed-cl-distributions, def:graph-distribution}
	For the Pauli type graph $G^{\mathrm{pauli}}$ and the family $L$ of Definition~\ref{def:pauli-question-distribution}, the Pauli question distribution equals $\mu^{G^{\mathrm{pauli}}}_{L,L}$.
\end{lemma}
\begin{proof}
	\uses{def:typed-cl-distributions, def:graph-distribution, def:cl-dist}
	Both distributions first use the uniform law on ordered pairs whose unordered pair is an edge of $G^{\mathrm{pauli}}$, and then push forward a common uniform seed in $V^{\mathrm{pauli}}$ by the same pair of maps.
\end{proof}

\begin{lemma}[The Pauli question distribution is a typed CL distribution]\label{lem:pauli-question-typed-cl}
	\lean{MIPStarRE.QPBT.pauliQuestionDistribution_isTypedCL}
	\uses{def:pauli-question-distribution, def:typed-cl-distributions}
	Fix an admissible tuple $\mathsf{qldparams} = (q,m,d)$ (Definition~\ref{def:admissible}) and let $L = \{L_t\}_{t \in \mathcal T^{\mathrm{pauli}}}$ be the family of Definition~\ref{def:pauli-question-distribution}. The family $L$ is a finite $\mathcal T^{\mathrm{pauli}}$-typed family of $3$-level conditionally linear functions, and the question distribution of $\mathfrak{G}^{\mathrm{Pauli}}_{\mathsf{qldparams}}$ equals the $(\mathcal T^{\mathrm{pauli}}, G^{\mathrm{pauli}})$-typed conditionally linear distribution $\mu^{G^{\mathrm{pauli}}}_{L,\,L}$ of Definition~\ref{def:typed-cl-distributions}.
\end{lemma}
\begin{proof}
	\uses{lem:pauli-question-cl-family, lem:pauli-question-typed-equality}
	Combine the common-level family assertion of Lemma~\ref{lem:pauli-question-cl-family} with the distribution identity of Lemma~\ref{lem:pauli-question-typed-equality}.
\end{proof}

\begin{definition}[Win predicate of the Pauli basis test]\label{def:pauli-win-predicate}
	\uses{def:pauli-question-distribution, def:ld-win-predicate, def:ms-game, def:indicator-vector, def:admissible}
	Fix an admissible tuple $\mathsf{qldparams} = (q,m,d)$ (Definition~\ref{def:admissible}) and let $M = 2^m$. Answers in the game $\mathfrak{G}^{\mathrm{Pauli}}_{\mathsf{qldparams}}$ take the following forms, depending on the question type; answers not of the prescribed form are rejected.
	% The paper's format table writes the line contents loosely as (u_W, s) and (u_W, s, v);
	% by the sampling procedure the actual contents are the canonical pairs (u_0, s) and
	% triples (u_0, s, v') with v' = pi_{i-1}(v). We state the canonical form.
	\begin{itemize}
		\item $(\mathsf{Point},W)$: content $y \in \F_q^m$; answer an element of $\F_q$.
		\item $(\mathsf{ALine},W)$: content $(u_0, s) \in \F_q^m \times \F_q$; answer a univariate polynomial $f\colon \F_q \to \F_q$ of degree at most $d$.
		\item $(\mathsf{DLine},W)$: content $(u_0, s, v') \in \F_q^m \times \F_q \times \F_q^m$; answer a univariate polynomial $f\colon \F_q \to \F_q$ of degree at most $md$.
		\item $\mathsf{Pair}$: content $\omega = (u_X, u_Z, r_X, r_Z) \in (\F_q^m)^2 \times \F_q^2$; answer $(\beta_X, \beta_Z) \in \F_2^2$.
		\item $(\mathsf{Pair},W)$: content $\omega$ as above; answer an element of $\F_2$.
		\item $\mathsf{Constraint}_i$: content $\omega$ as above; answer $(\alpha_{v_1}, \alpha_{v_2}, \alpha_{v_3}) \in \F_2^3$, indexed by the variable indices $(v_1,v_2,v_3)$ of $\mathsf{Constraint}_i$.
		\item $\mathsf{Variable}_j$: content $\omega$ as above; answer an element of $\F_2$.
		\item $(\mathsf{Pauli},W)$: content $0$; answer an element of $\F_q^M$.
	\end{itemize}
	The win predicate $D^{\mathrm{pauli}} = D^{\mathrm{pauli}}_{\mathsf{qldparams}}$ is the function of the tuple $(t_{\mathrm A}, x_{\mathrm A}, t_{\mathrm B}, x_{\mathrm B}, a_{\mathrm A}, a_{\mathrm B})$ determined by the following clauses, evaluated for $w \in \{\mathrm A, \mathrm B\}$ and $W \in \{X,Z\}$; in all cases where no action is indicated, the predicate accepts. In the last four clauses, the content of the player whose type lies in $\mathcal T^{\mathrm{ms}} \cup \{\mathsf{Pair}, (\mathsf{Pair},X), (\mathsf{Pair},Z)\}$ is a tuple $(u_X, u_Z, r_X, r_Z)$, from which the predicate first computes
	\begin{equation}
		\label{eq:gamma-value}
		\gamma = \tr\bigl( (\mathrm{ind}_m(u_X) \cdot r_X) \cdot (\mathrm{ind}_m(u_Z) \cdot r_Z) \bigr) \in \F_2\;,
	\end{equation}
	where $\mathrm{ind}_m(\cdot) \in \F_q^M$ is the indicator vector of the low-degree encoding (Chapter~\ref{chap:qpbt-algebra}), the inner dot denotes the dot product on $\F_q^M$, and $\tr\colon \F_q \to \F_2$ is the finite-field trace.
	\begin{enumerate}
		\item (\textbf{Consistency check}) If $t_{\mathrm A} = t_{\mathrm B}$, accept if and only if $a_{\mathrm A} = a_{\mathrm B}$.
		\item (\textbf{Low-degree check}) Let $\mathsf{ldparams} = (q,m,d,1)$; admissibility of $\mathsf{qldparams}$ gives $d \ge 1$ and $m \mid q$, so this is a parameter tuple as in Definition~\ref{def:ld-game}.
		\begin{enumerate}
			\item If $t_w = (\mathsf{Point},W)$ and $t_{\overline w} = (\mathsf{ALine},W)$, accept if $D^{\mathrm{ld}}_{\mathsf{ldparams}}$ accepts $(\mathsf{Point}, x_w, \mathsf{ALine}, x_{\overline w}, a_w, a_{\overline w})$.
			\item If $t_w = (\mathsf{Point},W)$ and $t_{\overline w} = (\mathsf{DLine},W)$, accept if $D^{\mathrm{ld}}_{\mathsf{ldparams}}$ accepts $(\mathsf{Point}, x_w, \mathsf{DLine}, x_{\overline w}, a_w, a_{\overline w})$.
		\end{enumerate}
		\item (\textbf{Consistency check}) If $t_w = (\mathsf{Point},W)$ and $t_{\overline w} = (\mathsf{Pauli},W)$, accept if $g_{a_{\overline w}}(x_w) = a_w$, where $g_{a_{\overline w}}$ is the low-degree encoding of $a_{\overline w} \in \F_q^M$ (Chapter~\ref{chap:qpbt-algebra}).
		\item (\textbf{Commutation check}) If $t_w = (\mathsf{Pair},W)$ and $t_{\overline w} = \mathsf{Pair}$, accept if $a_w = \beta_W$ or $\gamma \neq 0$.
		\item (\textbf{Consistency check}) If $t_w = (\mathsf{Point},W)$ and $t_{\overline w} = (\mathsf{Pair},W)$, accept if $\tr(a_w r_W) = a_{\overline w}$ or $\gamma \neq 0$.
		\item (\textbf{Magic Square check}) If $t_w = \mathsf{Constraint}_i$ and $t_{\overline w} = \mathsf{Variable}_j$, accept if $\gamma = 0$, or if $a_w$ satisfies the equation of $\mathsf{Constraint}_i$ and $\alpha_j = a_{\overline w}$.
		\item (\textbf{Consistency check}) If $t_w = (\mathsf{Point},W)$ and $t_{\overline w} = \mathsf{Variable}_j$, accept if $\gamma = 0$, or if $j = 1$, $W = X$, and $\tr(a_w r_X) = a_{\overline w}$, or if $j = 5$, $W = Z$, and $\tr(a_w r_Z) = a_{\overline w}$.
	\end{enumerate}
	The game $\mathfrak{G}^{\mathrm{Pauli}}_{\mathsf{qldparams}}$ is the typed game with type set $\mathcal T^{\mathrm{pauli}}$, the question distribution of Definition~\ref{def:pauli-question-distribution}, and win predicate $D^{\mathrm{pauli}}_{\mathsf{qldparams}}$.
\end{definition}

Note the one-sided gating on $\gamma$: the commutation check and the $\mathsf{Pair}$ consistency check accept automatically when $\gamma \neq 0$, while the Magic Square check and the $\mathsf{Variable}$ consistency check accept automatically when $\gamma = 0$.

\begin{lemma}[Completeness of the Pauli basis test]\label{lem:pauli-completeness}
	\lean{MIPStarRE.QPBT.pauliBasisTestSymm, MIPStarRE.QPBT.exists_spcc_value_one}
	\leanok
	\uses{def:pauli-question-distribution, def:pauli-win-predicate, def:admissible, def:spcc, def:tensor-product-value}
	For every admissible tuple $\mathsf{qldparams} = (q,m,d)$ (Definition~\ref{def:admissible}), the Pauli basis test $\mathfrak{G}^{\mathrm{Pauli}}_{\mathsf{qldparams}}$ has a value-$1$ SPCC strategy.
\end{lemma}
\begin{proof}
	\uses{thm:ms-from-ac, def:pauli-win-predicate, def:bracket, def:consistent-measurement, def:EPR, def:indicator-vector, lem:pauli-observable-expansion, lem:twisted-commutation}
	The strategy uses the state $\ket{\psi} = \ket{\mathrm{EPR}_q}^{\ot M} \ot \ket{\mathrm{EPR}_2}$ with $M = 2^m$. On a question $((\mathsf{Point},W), y)$ a player measures the generalized Pauli basis POVM $\{\tau^W_h\}_{h \in \F_q^M}$ and reports $g_h(y)$; on $((\mathsf{ALine},W), \ell)$ or $((\mathsf{DLine},W), \ell)$ they report the restriction of $g_h$ to the line $\ell$, expressed as a univariate polynomial through the canonical parameterization of $\ell$; and on $(\mathsf{Pauli},W)$ they report $h$ itself --- in the bracket notation of Definition~\ref{def:bracket}, these measurements are $\tau^W_{[g_\cdot(y) = a]} \ot I$, $\tau^W_{[(g_\cdot)|_\ell = f]} \ot I$, and $\tau^W_a \ot I$. For a content $\omega = (u_X, u_Z, r_X, r_Z)$, consider the two binary measurements $A_b = \tau^X_{[\tr(g_\cdot(u_X)\, r_X) = b]} \ot I$ and $B_b = \tau^Z_{[\tr(g_\cdot(u_Z)\, r_Z) = b]} \ot I$. Since $g_h(u) = h \cdot \mathrm{ind}_m(u)$ (Chapter~\ref{chap:qpbt-algebra}), the relation between the Pauli projections and observables gives $\mathcal O_A = A_0 - A_1 = \tau^X(\mathrm{ind}_m(u_X)\, r_X) \ot I$ and $\mathcal O_B = \tau^Z(\mathrm{ind}_m(u_Z)\, r_Z) \ot I$, and the twisted commutation relation yields
	\[
		\mathcal O_A \mathcal O_B = (-1)^\gamma\, \mathcal O_B \mathcal O_A
	\]
	with $\gamma$ as in Equation~\eqref{eq:gamma-value}: the two measurements commute when $\gamma = 0$ and anticommute when $\gamma = 1$. If $\gamma = 0$, define $E^{\mathsf{Pair},\omega}_{\beta_X,\beta_Z} = A_{\beta_X} B_{\beta_Z}$ and $E^{(\mathsf{Pair},X),\omega}_a = A_a$, $E^{(\mathsf{Pair},Z),\omega}_a = B_a$, and let the $\mathsf{Constraint}$ and $\mathsf{Variable}$ measurements be trivial (identity on the all-zeros outcome). If $\gamma = 1$, let the $\mathsf{Constraint}$ and $\mathsf{Variable}$ measurements be those produced by Theorem~\ref{thm:ms-from-ac} from the anticommuting consistent pair $(A, B)$, acting on the auxiliary $\ket{\mathrm{EPR}_2}$ factor, and let the $\mathsf{Pair}$ measurements be trivial. The strategy is symmetric and projective by construction; every measurement is either a post-processed Pauli basis measurement, a measurement supplied by Theorem~\ref{thm:ms-from-ac}, or (in the $\gamma = 0$ case) a product of two commuting consistent measurements, so the strategy is consistent, and commutativity along every edge of $G^{\mathrm{pauli}}$ is checked case by case --- for instance, $(\mathsf{Point},X)$ and $\mathsf{Variable}_1$ commute because both are $X$-basis measurements. Consistency and projectivity give the consistency checks with certainty, and the low-degree checks pass because the answers agree with the strategy determined by the polynomial $g_h$ in the classical game. For the remaining checks: when $\gamma = 0$ the commutation check passes by construction, and when $\gamma = 1$ the Magic Square check passes by Theorem~\ref{thm:ms-from-ac}; the consistency checks 5 and 7 pass with certainty because $\tau^W_{[\tr(g_\cdot(y)\, r_W) = \cdot]} \ot I$ coincides with $E^{(\mathsf{Pair},W),\omega}$, respectively with the embedded $\mathsf{Variable}_1$ and $\mathsf{Variable}_5$ measurements of Theorem~\ref{thm:ms-from-ac} (item~3 there). Hence the strategy has value $1$.
\end{proof}

We now state the main soundness property of the Pauli basis test; the statement is an adaptation of the self-testing theorem of~\cite{Natarajan2018NEEXP} (Theorem~6.4 there). Unlike Theorem~\ref{lem:ld-soundness}, no projectivity assumption is made on the strategy.

\begin{theorem}[Soundness of the Pauli basis test]\label{thm:pauli}
	\uses{def:pauli-question-distribution, def:pauli-win-predicate, def:admissible, def:tensor-product-value, def:EPR, lem:pauli-observable-expansion, def:povm-distance}
	There exists a function
	\[
		\delta_{\mathrm{qld}}(\eps, m, d, q) = a\,(md)^a \bigl( \eps^b + q^{-b} + 2^{-bmd} \bigr)
	\]
	for universal constants $a \ge 1$ and $0 < b < 1$ such that the following holds. For all admissible parameter tuples $\mathsf{qldparams} = (q,m,d)$ and for all strategies $\mathcal S = (\ket\psi, A, B)$ for the game $\mathfrak{G}^{\mathrm{Pauli}}_{\mathsf{qldparams}}$ that succeed with probability at least $1 - \eps$, there exist local isometries $\phi_{\mathrm A}\colon \hilb_{\mathrm A} \to \hilb_{\mathrm A'} \ot \hilb_{\mathrm A''}$ and $\phi_{\mathrm B}\colon \hilb_{\mathrm B} \to \hilb_{\mathrm B'} \ot \hilb_{\mathrm B''}$, where $\ket\psi \in \hilb_{\mathrm A} \ot \hilb_{\mathrm B}$ and $\hilb_{\mathrm A''}, \hilb_{\mathrm B''} \cong (\C^q)^{\ot M}$ with $M = 2^m$, and a state $\ket{\mathrm{aux}} \in \hilb_{\mathrm A'} \ot \hilb_{\mathrm B'}$, such that:
	\begin{enumerate}
		\item $\bigl\| \phi_{\mathrm A} \ot \phi_{\mathrm B} \ket\psi - \ket{\mathrm{aux}} \ot \ket{\mathrm{EPR}_q}^{\ot M} \bigr\| \le \delta_{\mathrm{qld}}(\eps, m, d, q)$;
		% The paper here has the typo "phi_alice \ot phi_B"; both maps are the isometries above.
		\item letting $\tilde A^x_a = \phi_{\mathrm A} A^x_a \phi_{\mathrm A}^\dagger$ and $\tilde B^y_b = \phi_{\mathrm B} B^y_b \phi_{\mathrm B}^\dagger$, for $W \in \{X, Z\}$,
		\begin{align*}
			\tilde A^{(\mathsf{Pauli},W)}_u \ot I_{\mathrm{B'B''}} &\approx_{\delta_{\mathrm{qld}}} (\tau^W_u)_{\mathrm A''} \ot I_{\mathrm{A'B'B''}}\;,\\
			I_{\mathrm{A'A''}} \ot \tilde B^{(\mathsf{Pauli},W)}_u &\approx_{\delta_{\mathrm{qld}}} I_{\mathrm{A'A''B'}} \ot (\tau^W_u)_{\mathrm B''}\;,
		\end{align*}
		where the $\approx_{\delta_{\mathrm{qld}}}$ statements hold with respect to the state $\ket{\mathrm{aux}}_{\mathrm{A'B'}} \ot \ket{\mathrm{EPR}_q}^{\ot M}_{\mathrm{A''B''}}$ and the answer summation is over $u \in \F_q^M$.
	\end{enumerate}
\end{theorem}

The proof of Theorem~\ref{thm:pauli} is the adaptation, to low individual degree, of the analysis of the quantum low-degree test of Natarajan and Vidick; it is developed in Chapters~\ref{chap:qpbt-observables} and~\ref{chap:qpbt-combining} and concluded in Chapter~\ref{chap:qpbt-extraction}, where the argument is complete only modulo the unresolved divisibility obligation documented in \cite{gap:qpbt_ld-dimension-divisibility}. It relies on the classical soundness theorem (Theorem~\ref{lem:ld-soundness}) and on the rigidity of the Magic Square game (Theorem~\ref{thm:ms-rigidity}).

The strategies appearing in Lemma~\ref{lem:pauli-completeness} and in the conclusion of Theorem~\ref{thm:pauli} are phrased in terms of qudit Pauli measurements and maximally entangled states of $q$-dimensional qudits. For the intended application it is more convenient to work with qubits: since the field sizes $q$ used here are powers of $2$, the conclusions can be restated in terms of qubit EPR pairs and qubit Pauli measurements.

\begin{corollary}[Qubit form of the soundness of the Pauli basis test]\label{cor:pauli-binary}
	\lean{MIPStarRE.QPBT.pauli_soundness_qubit}
	\leanok
	\uses{def:pauli-question-distribution, def:pauli-win-predicate, def:admissible, def:tensor-product-value, def:EPR, def:povm-distance, def:binary-representation}
	There exists a function
	\[
		\delta_{\mathrm{qld}}(\eps, m, d, q) = a\,(md)^a \bigl( \eps^b + q^{-b} + 2^{-bmd} \bigr)
	\]
	% The paper's displayed function here writes the argument order (eps, q, m, d); the
	% formula is the same, and its own item 1 reverts to (eps, m, d, q). See the remark below.
	for universal constants $a \ge 1$ and $0 < b < 1$ such that the following holds. For all admissible parameter tuples $\mathsf{qldparams} = (q,m,d)$ and for all strategies $\mathcal S = (\ket\psi, A, B)$ for the game $\mathfrak{G}^{\mathrm{Pauli}}_{\mathsf{qldparams}}$ that succeed with probability at least $1 - \eps$, there exist local isometries $\phi_{\mathrm A}\colon \hilb_{\mathrm A} \to \hilb_{\mathrm A'} \ot \hilb_{\mathrm A''}$ and $\phi_{\mathrm B}\colon \hilb_{\mathrm B} \to \hilb_{\mathrm B'} \ot \hilb_{\mathrm B''}$, where $\ket\psi \in \hilb_{\mathrm A} \ot \hilb_{\mathrm B}$ and $\hilb_{\mathrm A''}, \hilb_{\mathrm B''} \cong (\C^2)^{\ot M \log q}$ with $M = 2^m$, and a state $\ket{\mathrm{aux}} \in \hilb_{\mathrm A'} \ot \hilb_{\mathrm B'}$, such that:
	\begin{enumerate}
		\item $\bigl\| \phi_{\mathrm A} \ot \phi_{\mathrm B} \ket\psi - \ket{\mathrm{aux}} \ot \ket{\mathrm{EPR}_2}^{\ot M \log q} \bigr\| \le \delta_{\mathrm{qld}}(\eps, m, d, q)$;
		\item letting $\tilde A^x_a = \phi_{\mathrm A} A^x_a \phi_{\mathrm A}^\dagger$ and $\tilde B^y_b = \phi_{\mathrm B} B^y_b \phi_{\mathrm B}^\dagger$, for $W \in \{X, Z\}$,
		\begin{align*}
			\tilde A^{(\mathsf{Pauli},W)}_u \ot I_{\mathrm{B'B''}} &\approx_{\delta_{\mathrm{qld}}} (\sigma^W_u)_{\mathrm A''} \ot I_{\mathrm{A'B'B''}}\;,\\
			I_{\mathrm{A'A''}} \ot \tilde B^{(\mathsf{Pauli},W)}_u &\approx_{\delta_{\mathrm{qld}}} I_{\mathrm{A'A''B'}} \ot (\sigma^W_u)_{\mathrm B''}\;,
		\end{align*}
		where the $\approx_{\delta_{\mathrm{qld}}}$ statements hold with respect to the state $\ket{\mathrm{aux}}_{\mathrm{A'B'}} \ot \ket{\mathrm{EPR}_2}^{\ot M \log q}_{\mathrm{A''B''}}$ and the answer summation is over $u \in \F_q^M$. Here $\sigma^W_u$ denotes the tensor product of single-qubit Pauli measurements
		\[
			\bigotimes_{i=1}^{M} \bigotimes_{j=1}^{\log q} \sigma^W_{u_{ij}}\;,
		\]
		where $\kappa(u_i) = (u_{ij})_{1 \le j \le \log q}$ under the bijection $\kappa\colon \F_q \to \F_2^{\log q}$ of Chapter~\ref{chap:qpbt-algebra}.
	\end{enumerate}
\end{corollary}
\begin{proof}
	\uses{thm:pauli, lem:pauli-binary}
	By Lemma~\ref{lem:pauli-binary}, since $q$ is a power of $2$ there are local isometries that map $\ket{\mathrm{EPR}_q}^{\ot M}$ to $\ket{\mathrm{EPR}_2}^{\ot M \log q}$ and map the generalized Pauli measurements $\tau^W_u$ to the qubit tensor products $\bigotimes_{i=1}^M \bigotimes_{j=1}^{\log q} \sigma^W_{u_{ij}}$, where $\kappa(u_i) = (u_{i1},\ldots,u_{i \log q})$. Composing these isometries with those produced by Theorem~\ref{thm:pauli} gives the statement.
\end{proof}

\begin{remark}[Notational drift in the soundness functions]\label{rem:delta-qld-argument-order}
	The paper writes the soundness function of the Pauli basis test with two different argument orders: $\delta_{\mathrm{qld}}(\eps, m, d, q)$ in Theorem~\ref{thm:pauli} and Lemma~\ref{lem:delta-bound}, and $\delta_{\mathrm{qld}}(\eps, q, m, d)$ in the displayed function of Corollary~\ref{cor:pauli-binary}, whose own item~1 reverts to $\delta_{\mathrm{qld}}(\eps, m, d, q)$; the formula is the same throughout, and here the order $(\eps, m, d, q)$ is used everywhere. Note also two deliberate asymmetries with the classical soundness function of Theorem~\ref{lem:ld-soundness}: there the polynomial prefactor is $a(dmk)^a$ and the exponent satisfies $0 < b \le 1$, while in Theorems~\ref{thm:pauli} and Corollary~\ref{cor:pauli-binary} the prefactor is $a(md)^a$ and the exponent satisfies the strict inequality $0 < b < 1$.
\end{remark}

\section{Canonical parameters of the Pauli basis test}
% Paper label sec:qld-complexity kept for cross-reference fidelity (the complexity
% statements of that subsection are outside the scope of this development).
\label{sec:qld-complexity}

We specify a canonical setting of the parameter tuple $\mathsf{qldparams}$ as a function of an integer $R$, and bound the resulting soundness function. All logarithms are base $2$.

\begin{definition}[Canonical parameters of the Pauli basis test]\label{def:introparams}
	\lean{MIPStarRE.QPBT.introParamsC, MIPStarRE.QPBT.introParamsTuple, MIPStarRE.QPBT.AdmissibleParams.ofTuple, MIPStarRE.QPBT.introParams}
	\leanok
	\uses{thm:pauli}
	Let $a, b$ be the universal constants of Theorem~\ref{thm:pauli}, and let $c$ be the smallest \emph{even} integer that is at least $(b+a)/b$. For every integer $R \ge 4$, define the tuple $\mathsf{introparams}(R) = (q, m, d)$ by
	\begin{itemize}
		\item $q = 2^k$ for $k = c \lceil \log \log R \rceil + 1$,
		\item $m = 2^j$, where $j$ is the unique integer with $2^j \le c \lceil \log R \rceil + 1 < 2^{j+1}$,
		\item $d = 1$.
	\end{itemize}
	The letter $k$ here denotes the exponent of the field size and is unrelated to the simultaneity parameter $k$ of the classical game. With this choice the Pauli basis test certifies the presence of an $M = 2^m$-qudit maximally entangled state of local dimension $q$, equivalently an $(M \log q)$-qubit EPR state, and the choice $d = 1$ runs the low individual degree test multilinearly.
\end{definition}

\begin{lemma}[Soundness bound for the canonical parameters]\label{lem:delta-bound}
	\lean{MIPStarRE.QPBT.introParamsTuple_isAdmissible, MIPStarRE.QPBT.le_two_pow_introParams_m, MIPStarRE.QPBT.exists_deltaQld_introParams_bound}
	\leanok
	\uses{def:introparams, def:admissible, thm:pauli}
	For all integers $R \ge 4$, the parameter tuple $\mathsf{introparams}(R) = (q,m,d)$ is admissible (Definition~\ref{def:admissible}) and satisfies $2^m \ge R$. Furthermore, there exist universal constants $a' \ge 1$ and $0 < b' \le 1$ such that, for every $\eps \ge 0$, the function $\delta_{\mathrm{qld}}(\eps, m, d, q)$ of Theorem~\ref{thm:pauli} is at most
	\[
		a' \bigl( (\log R)^{a'} \cdot \eps^{b'} + (\log R)^{-b'} \bigr)\;.
	\]
\end{lemma}
\begin{proof}
	\uses{def:introparams, def:admissible, def:admissible-size, thm:pauli}
	Since $c$ is even, the exponent $k = c \lceil \log\log R \rceil + 1$ is odd, so $q = 2^k$ is an admissible field size; $m$ divides $q$ because $2^j \le c \lceil \log R \rceil + 1 \le 2 \log^c R \le 2^k$, so $j \le k$ and both are powers of $2$; and $d = 1 \ge 1$. Hence the tuple is admissible. Since $m = 2^j \ge (c/2) \log R$, we get $2^m \ge R^{c/2} \ge R$, using $c \ge 2$.

	For the soundness bound, write $\delta_{\mathrm{qld}}(\eps,m,d,q) = a(md)^a \eps^b + a(md)^a q^{-b} + a(md)^a 2^{-bmd}$ and use $d = 1$, $m \le 2c \log R$, and $q = 2^k \ge \log^c R$. The first term is at most $a(2c)^a (\log R)^a \eps^b$. The second term is at most $a(2c)^a (\log R)^{a - cb} \le a(2c)^a (\log R)^{-b}$, since $cb - a \ge b$ by the choice of $c$. For the third term, $m \ge (c/2) \log R \ge \log R$ gives $2^{-bm} \le R^{-b}$, so
	\[
		a(md)^a\, 2^{-bm}
		\le a(2c)^a (\log R)^a R^{-b}
		= a(2c)^a \bigl( (\log R)^{a+b} R^{-b} \bigr) (\log R)^{-b}
		\le a(2c)^a\, C\, (\log R)^{-b}\;,
	\]
	where $C = \sup_{R' \ge 4} (\log R')^{a+b} (R')^{-b}$ is finite --- the function is continuous in $R'$ and tends to $0$ as $R' \to \infty$ --- and depends only on the universal constants $a$ and $b$. Summing the three bounds, the statement follows with $a' = a(2c)^a (1 + C)$ and $b' = b$.
\end{proof}

\begin{remark}[A comparison of exponents in the source proof of Lemma~\ref{lem:delta-bound}]\label{rem:delta-bound-exponent-comparison}
	The source proof bounds the third term of $\delta_{\mathrm{qld}}$ by asserting $2^{-bmd} \le q^{-b}$ ``because $k \le m$'', comparing the field exponent $k = c \lceil \log\log R \rceil + 1$ with $m = 2^j$. The inequality $k \le m$ does not follow from Definition~\ref{def:introparams}: if the universal constants of Theorem~\ref{thm:pauli} were $a = 1$ and $b = 1/2$, then $c = 4$, and $R = 5$ gives $k = 9$ and $m = 8$. The proof above instead bounds $2^{-bm} \le R^{-b}$ directly from $m \ge (c/2) \log R \ge \log R$, which yields the same conclusion after enlarging the universal constant $a'$.  The comparison and correction are documented in \cite{gap:qpbt_delta-bound-exponent-comparison}.
\end{remark}

\section{Formalization support for finite sampling}

The following finite-distribution identities are not named statements of the
source article.  They make explicit the push-forward calculations used in the
low-degree question and line-point laws above.

\begin{lemma}[Extensionality of finite distributions]\label{lem:distribution-ext-support}
	\lean{MIPStarRE.QPBT.Distribution.ext_of_support_of_weight}
	\leanok
	Two finite distributions with the same support and the same weight at every
	point are equal.
\end{lemma}
\begin{proof}
	\leanok
	The two laws have the same finite support and assign the same mass to every
	point, so they coincide.
\end{proof}

\begin{lemma}[Composition of push-forwards]\label{lem:distribution-map-comp}
	\lean{MIPStarRE.QPBT.Distribution.map_map}
	\leanok
	Let $\mu$ be a finite distribution on $A$, and let $e\colon A\to B$ and
	$f\colon B\to C$. Then
	\[
		f_\ast(e_\ast\mu)=(f\circ e)_\ast\mu.
	\]
\end{lemma}
\begin{proof}
	\leanok
	\uses{lem:distribution-ext-support}
	The supports agree by composition of finite images.  For each $c\in C$,
	partition the fiber of $f\circ e$ over $c$ by the intermediate value $e(a)$;
	the corresponding sums of weights agree.
\end{proof}

\begin{lemma}[Weights of the uniform law]\label{lem:uniform-distribution-weight}
	\lean{MIPStarRE.QPBT.uniformDistribution_weight_apply}
	\leanok
	If $A$ is a nonempty finite set, then its uniform distribution assigns weight
	$1/|A|$ to every $a\in A$.
\end{lemma}
\begin{proof}
	\leanok
	This is the defining weight of the uniform distribution on a nonempty finite
	set.
\end{proof}

\begin{lemma}[Uniform push-forward along a balanced map]\label{lem:uniform-map-equal-fibers}
	\lean{MIPStarRE.QPBT.uniformDistribution_map_of_card_fiber}
	\leanok
	Let $A$ and $B$ be nonempty finite sets and let $e\colon A\to B$. If there is
	an integer $c$ such that $|e^{-1}(b)|=c$ for every $b\in B$, then
	$e_\ast\operatorname{Unif}(A)=\operatorname{Unif}(B)$.
\end{lemma}
\begin{proof}
	\leanok
	\uses{lem:distribution-ext-support, lem:uniform-distribution-weight}
	Every fiber is nonempty, so $e$ is surjective and both supports are $B$.
	Counting the fibers gives $|A|=c|B|$; summing the weight $1/|A|$ over a
	$c$-element fiber gives $1/|B|$.
\end{proof}

\begin{lemma}[Uniformity under an equivalence]\label{lem:uniform-map-equivalence}
	\lean{MIPStarRE.QPBT.uniformDistribution_map_equiv}
	\leanok
	An equivalence $E\colon A\simeq B$ between nonempty finite sets satisfies
	$E_\ast\operatorname{Unif}(A)=\operatorname{Unif}(B)$.
\end{lemma}
\begin{proof}
	\leanok
	\uses{lem:uniform-map-equal-fibers}
	Every fiber of an equivalence has cardinality one, so the result follows from
	Lemma~\ref{lem:uniform-map-equal-fibers}.
\end{proof}

\begin{lemma}[First marginal of a uniform product]\label{lem:uniform-product-first-marginal}
	\lean{MIPStarRE.QPBT.uniformDistribution_map_fst}
	\leanok
	For nonempty finite sets $A$ and $B$, the first marginal of
	$\operatorname{Unif}(A\times B)$ is $\operatorname{Unif}(A)$.
\end{lemma}
\begin{proof}
	\leanok
	\uses{lem:uniform-map-equal-fibers}
	Every fiber of the first projection has cardinality $|B|$; apply
	Lemma~\ref{lem:uniform-map-equal-fibers}.
\end{proof}

\begin{lemma}[Second marginal of a uniform product]\label{lem:uniform-product-second-marginal}
	\lean{MIPStarRE.QPBT.uniformDistribution_map_snd}
	\leanok
	For nonempty finite sets $A$ and $B$, the second marginal of
	$\operatorname{Unif}(A\times B)$ is $\operatorname{Unif}(B)$.
\end{lemma}
\begin{proof}
	\leanok
	\uses{lem:uniform-map-equal-fibers}
	Every fiber of the second projection has cardinality $|A|$; apply
	Lemma~\ref{lem:uniform-map-equal-fibers}.
\end{proof}

\begin{lemma}[Dependent uniform sampling as a product push-forward]\label{lem:uniform-bind-map}
	\lean{MIPStarRE.QPBT.bind_uniformDistribution_map}
	\leanok
	Let $A$ and $B$ be nonempty finite sets and let $g\colon A\to B\to C$.
	Sampling $a$ uniformly from $A$, then $b$ uniformly from $B$, and returning
	$g(a,b)$ has the same law as pushing $\operatorname{Unif}(A\times B)$ forward
	along $(a,b)\mapsto g(a,b)$.
\end{lemma}
\begin{proof}
	\leanok
	\uses{lem:distribution-ext-support, lem:uniform-distribution-weight}
	Both supports are the image of $A\times B$ under $g$.  At each $c\in C$,
	both weights are the number of pairs $(a,b)$ with $g(a,b)=c$, divided by
	$|A||B|$.
\end{proof}

\begin{lemma}[Uniform first marginal after an equivalence]
	\label{lem:uniform-first-marginal-after-equivalence}
	\lean{MIPStarRE.QPBT.uniformDistribution_map_fst_of_equiv}
	\leanok
	Let $A$, $B$, and $C$ be nonempty finite sets, let $E\colon A\simeq B\times C$,
	and suppose $f(a)$ is the first component of $E(a)$ for every $a\in A$. Then
	$f_\ast\operatorname{Unif}(A)=\operatorname{Unif}(B)$.
\end{lemma}
\begin{proof}
	\leanok
	\uses{lem:distribution-map-comp, lem:uniform-map-equivalence, lem:uniform-product-first-marginal}
	Push the uniform law along $E$, compose with the first projection, and apply
	uniformity under the equivalence followed by the first-marginal identity.
\end{proof}

\begin{lemma}[Exact fibers of the coordinate-index map]\label{lem:chi-seed-fibers}
	\lean{MIPStarRE.QPBT.LdParams.seedFiberCard_pos, MIPStarRE.QPBT.LdParams.seedFiberProduct_eq}
	\lean{MIPStarRE.QPBT.seedFiberEquiv, MIPStarRE.QPBT.seedFiberEquiv_fst}
	\lean{MIPStarRE.QPBT.seedOfIndexResidue, MIPStarRE.QPBT.seedFiberEquiv_seedOfIndexResidue, MIPStarRE.QPBT.chiIndex_seedOfIndexResidue}
	\leanok
	\uses{def:ld-game, def:ld-question-distribution, def:binary-representation}
	Write $\chi_0(s)=\chi(s)-1$ for the corresponding zero-based coordinate
	index. Every scalar seed has a unique
	representation as a pair
	\[
		(i,r)\in\{0,\ldots,m-1\}\times\{0,\ldots,q/m-1\},
	\]
	and the first component is $\chi_0(s)$. Conversely, every such pair determines
	a unique scalar seed. In particular, the residue set is nonempty and
	$q=m(q/m)$.
\end{lemma}
\begin{proof}
	\leanok
	\uses{def:binary-representation}
	Use the fixed bijection between $\F_q$ and $\{0,\ldots,q-1\}$, the equality
	$q=m(q/m)$, and quotient--remainder decomposition with divisor $q/m$.
\end{proof}

\begin{lemma}[Uniformity of the coordinate index]\label{lem:chi-index-uniform}
	\lean{MIPStarRE.QPBT.uniformDistribution_map_chiIndex}
	\leanok
	\uses{def:ld-question-distribution}
	For a uniformly random $s\in\F_q$, the zero-based index $\chi_0(s)$ is
	uniformly distributed on $\{0,\ldots,m-1\}$.
\end{lemma}
\begin{proof}
	\leanok
	\uses{lem:chi-seed-fibers, lem:uniform-first-marginal-after-equivalence}
	Under the equivalence of Lemma~\ref{lem:chi-seed-fibers}, $\chi_0$ is the first
	projection. Apply Lemma~\ref{lem:uniform-first-marginal-after-equivalence}.
\end{proof}

\begin{lemma}[Idempotence of the prefix projection]
	\label{lem:prefix-projection-idempotent}
	\lean{MIPStarRE.QPBT.prefixProjection_idempotent}
	\leanok
	\uses{def:ld-question-distribution}
	For every coordinate $i$ and direction $v$,
	$\pi_i(\pi_i(v))=\pi_i(v)$.
\end{lemma}
\begin{proof}
	\leanok
	A coordinate before $i$ is set to zero by both sides, while every coordinate
	at or after $i$ is retained by both sides.
\end{proof}

\begin{lemma}[Idempotence of the canonical representative]
	\label{lem:line-representative-idempotent}
	\lean{MIPStarRE.QPBT.lineRepMap_apply_self}
	\leanok
	\uses{def:line-representative}
	For all $u,v\in\F_q^m$,
	$L^{\mathrm{Ln}}_v(L^{\mathrm{Ln}}_v(u))=L^{\mathrm{Ln}}_v(u)$.
\end{lemma}
\begin{proof}
	\leanok
	The canonical representative map is a projection onto the chosen complement
	of the span of $v$, and every projection is idempotent.
\end{proof}

\begin{definition}[Seed-bearing line descriptions]\label{def:ld-line-description}
	\lean{MIPStarRE.QPBT.LineKind, MIPStarRE.QPBT.LineDesc}
	\lean{MIPStarRE.QPBT.aLineDescOf, MIPStarRE.QPBT.dLineDescOf}
	\leanok
	\uses{def:ld-question-distribution, def:line-representative, lem:line-representative-idempotent}
	An axis-line description records a canonical base point and a scalar seed,
	whose coordinate index selects the axis direction. A diagonal-line
	description records a canonical base point, a scalar seed, and a direction
	whose coordinates preceding the selected index vanish. The axis and diagonal
	CL outputs determine such descriptions by taking their canonical
	representatives.
\end{definition}

\begin{lemma}[Prefix property of a diagonal description]
	\label{lem:line-description-prefix}
	\lean{MIPStarRE.QPBT.LineDesc.diagonal_prefix_zero}
	\leanok
	\uses{def:ld-line-description}
	If a seed-bearing line description is diagonal, then every coordinate of its
	direction preceding the index selected by its seed is zero.
\end{lemma}
\begin{proof}
	\leanok
	This is the prefix-zero property carried by a diagonal line description.
\end{proof}
